\section{Aggregation across consumers}\label{sec:05-aggregation}

The AIDS of \cref{sec:05-aids} describes the choices of a single consumer. Most data in industrial
organization are aggregates: the sales of a product in a city and a month, or national spending on a
group of goods in a year. This section asks when such aggregates behave as if one consumer had chosen
them. The answer decides two things. It decides whether the share equations of \eqref{eq:05-aids-share}
can be fitted to market data, and whether the welfare measures of \cref{sec:05-cv-ev}, computed from the
fitted system, describe anyone's welfare.

We index households by $h = 1, \dots, H$. Household $h$ has total expenditure $I_h$, demands $q_{ih}$
and budget shares $w_{ih} = \frac{p_i q_{ih}}{I_h}$. All households face the same price vector $p$.
Mean expenditure is $\bar I = \frac{1}{H}\sum_h I_h$. Hicksian demands keep the notation $h_i(p,u)$ of
\cref{sec:05-duality}, always written with their arguments, so they cannot be confused with the household
index.

\subsection{The aggregation problem}\label{sec:05-agg-problem}

Aggregate demand for good $i$ is the sum of household demands,
\begin{equation}\label{eq:05-agg-demand}
  Q_i(p, I_1, \dots, I_H) = \sum_{h=1}^{H} q_{ih}(p, I_h).
\end{equation}
In general it depends on the whole distribution of expenditure. To see why, transfer a small amount
$\dd t$ of expenditure from household $h'$ to household $h$, holding total expenditure fixed. Aggregate
demand changes by
\begin{equation*}
  \dd Q_i = \left(\frac{\partial q_{ih}}{\partial I_h} - \frac{\partial q_{ih'}}{\partial I_{h'}}\right)\dd t .
\end{equation*}
The change is zero only if the two households have the same marginal propensity to spend on good $i$
at their current budgets. Rich and poor households usually differ in this respect, so redistribution
moves market demand even when total spending does not change.

We want two things at once. Households should be allowed to differ in income and in demographics,
because they do. Yet we would like market data to look as if one consumer had generated them. A
\emph{representative consumer} is a fictional utility-maximizing consumer whose demand, evaluated at
current prices and at a scalar summary of the expenditure distribution, equals aggregate demand. We call
the representative consumer \emph{normative} when, in addition, its expenditure function measures the
welfare of the population, for example as the sum of household compensating variations. The positive
property licenses estimating a demand system on aggregate data. The normative property licenses welfare
analysis with the estimated system. The next subsections show how far each property can be pushed.

\subsection{The Gorman polar form and exact aggregation}\label{sec:05-gorman}

\begin{definition}[Gorman polar form]\label{def:05-gorman}
Households' preferences have the \emph{Gorman polar form} if their expenditure functions can be written
as
\begin{equation}\label{eq:05-gorman-form}
  e_h(p, u_h) = a_h(p) + b(p)\, u_h , \qquad h = 1, \dots, H.
\end{equation}
Here each $a_h(p)$ and the common function $b(p) > 0$ are homogeneous of degree one in $p$. On the
maintained domain of prices and utilities, each $e_h$ is increasing and concave in $p$. Preferences of
this form are called \emph{quasi-homothetic}.
\end{definition}

The function $a_h$ may differ across households, for example with family size or age. The function $b$
must be common. Under these conditions each $e_h$ is a valid expenditure function on the maintained
domain, and it is increasing in $u_h$ because $b > 0$. A convenient sufficient condition is that each
$a_h$ and $b$ is positive, increasing and concave, with $u_h \ge 0$. Then $e_h$ inherits monotonicity and
concavity in $p$ from $a_h$ and $b$. This subclass has $\frac{\partial \ln b}{\partial p_i} \ge 0$, which
by \eqref{eq:05-gorman-demand} below rules out inferior goods. The general form allows them on a local
domain. The name
reflects the relation to homotheticity. A homothetic expenditure function has the form $b(p)u$, so
budget shares do not depend on the budget. Adding the fixed cost $a_h(p)$ keeps Engel curves straight
but lets them miss the origin. The fixed cost is subsistence expenditure, the cost of living at $u_h = 0$
\citep[\S 5.4, pp.~142--144, and \S 6.1, eqs.~1.4--1.7, pp.~150--151]{DeatonMuellbauer1980b}.

\begin{proposition}[Exact aggregation, \citealp{Gorman1953,Gorman1961}]\label{prop:05-gorman}
Suppose every household has an interior optimum.
\begin{enumerate}[label=(\roman*)]
  \item If every household's expenditure function has the Gorman polar form \eqref{eq:05-gorman-form}
  with a common $b(p)$, aggregate demand depends on the distribution of expenditure only through
  $\sum_h I_h$. It is the demand of a representative consumer with expenditure function
  \begin{equation}\label{eq:05-gorman-rep}
    e(p, U) = \sum_{h=1}^{H} a_h(p) + b(p)\, U ,
  \end{equation}
  and the compensating variation of this consumer for any price change equals the sum of the
  households' compensating variations.
  \item Conversely, suppose aggregate demand depends only on $p$ and $\sum_h I_h$ for all distributions
  of expenditure in an open set. Then on that set household Engel curves are straight lines with slopes
  common to all households, and preferences have the Gorman polar form with a common $b(p)$. This $b$
  need not be increasing or concave.
\end{enumerate}
\end{proposition}

\begin{proof}
(i) Inverting $I_h = e_h(p, u_h)$ gives the indirect utility function
\begin{equation*}
  v_h(p, I_h) = \frac{I_h - a_h(p)}{b(p)} .
\end{equation*}
Roy's identity (\cref{prop:05-roy}) uses two partial derivatives:
\begin{align*}
  \frac{\partial v_h}{\partial I_h} &= \frac{1}{b(p)}, \\
  \frac{\partial v_h}{\partial p_i} &= -\frac{1}{b(p)}\frac{\partial a_h}{\partial p_i}
     - \frac{I_h - a_h(p)}{b(p)^2}\frac{\partial b}{\partial p_i} .
\end{align*}
Taking minus their ratio,
\begin{equation}\label{eq:05-gorman-demand}
  q_{ih}(p, I_h) = -\frac{\partial v_h}{\partial p_i}\left(\frac{\partial v_h}{\partial I_h}\right)^{-1}
  = \frac{\partial a_h}{\partial p_i} + \frac{\partial \ln b}{\partial p_i}\bigl(I_h - a_h(p)\bigr).
\end{equation}
The marginal propensity to spend, $\frac{\partial q_{ih}}{\partial I_h} = \frac{\partial \ln b}{\partial
p_i}$, depends on prices alone. It is the same for every household and at every budget. Summing
\eqref{eq:05-gorman-demand} over households gives
\begin{equation*}
  Q_i = \sum_h \frac{\partial a_h}{\partial p_i}
      + \frac{\partial \ln b}{\partial p_i}\Bigl(\sum_h I_h - \sum_h a_h(p)\Bigr),
\end{equation*}
which depends on budgets only through their sum.

Define $A(p) = \sum_h a_h(p)$ and $U = \sum_h u_h$. The function $e(p, U) = A(p) + b(p)U$ equals
$\sum_h e_h(p, u_h)$. It is therefore homogeneous of degree one, increasing and concave in $p$, and
increasing in $U$, so it is an expenditure function.
By Shephard's lemma (\cref{prop:05-shephard}) its Hicksian demand for good $i$ is $\frac{\partial
A}{\partial p_i} + \frac{\partial b}{\partial p_i}U$, the sum of the household Hicksian demands. Its
indirect utility at the total budget is $\frac{\sum_h I_h - A(p)}{b(p)} = \sum_h v_h(p, I_h)$, and Roy's
identity applied to it reproduces $Q_i$.

For a price change from $p^0$ to $p^1$, with initial utilities $u_h^0$, household $h$ has
\begin{equation*}
  CV_h = e_h(p^1, u_h^0) - e_h(p^0, u_h^0) = a_h(p^1) - a_h(p^0) + \bigl[b(p^1) - b(p^0)\bigr] u_h^0 .
\end{equation*}
Summing over households gives
\begin{equation*}
  \sum_h CV_h = A(p^1) - A(p^0) + \bigl[b(p^1) - b(p^0)\bigr] U^0 = e(p^1, U^0) - e(p^0, U^0),
\end{equation*}
the compensating variation of the representative consumer.

(ii) We sketch the argument. If $Q_i$ depends only on $\sum_h I_h$, every transfer between two
households leaves it unchanged. By the transfer argument of \cref{sec:05-agg-problem},
$\frac{\partial q_{ih}}{\partial I_h} = \frac{\partial q_{ih'}}{\partial I_{h'}}$ for every pair of
households and every pair of budgets in the open set. Fix the budget of $h'$ and vary that of $h$. The
derivative of $q_{ih}$ does not change, so it is a function $c_i(p)$ of prices alone, common to all
households. Integrating gives straight Engel curves $q_{ih} = d_{ih}(p) + c_i(p) I_h$. Recovering
preferences from these demands, with the symmetry and homogeneity restrictions of
\cref{prop:05-slutsky}, yields the Gorman polar form \citep{Gorman1961}.
\end{proof}

\begin{example}[The linear expenditure system]\label{ex:05-les}
The linear expenditure system of \citet{Stone1954} has expenditure function
\begin{equation*}
  e_h(p, u_h) = \sum_k p_k \underline{q}_{kh} + u_h \prod_k p_k^{m_k},
  \qquad m_k > 0, \quad \sum_k m_k = 1 ,
\end{equation*}
where $\underline{q}_{kh}$ is the subsistence quantity of good $k$ for household $h$. It has the Gorman
polar form with $a_h(p) = \sum_k p_k\underline{q}_{kh}$ and $b(p) = \prod_k p_k^{m_k}$. Since
$\frac{\partial \ln b}{\partial p_i} = \frac{m_i}{p_i}$, the demand function \eqref{eq:05-gorman-demand}
becomes
\begin{equation*}
  p_i q_{ih} = p_i \underline{q}_{ih} + m_i\Bigl(I_h - \sum_k p_k \underline{q}_{kh}\Bigr).
\end{equation*}
Each household first buys its subsistence bundle and then spends the fraction $m_i$ of its remaining
budget on good $i$. Households differ in subsistence quantities but share the marginal budget shares
$m_i$. Panel (a) of \cref{fig:05-gorman-piglog} draws two such Engel curves.
\end{example}

Exact aggregation comes at a price. Engel curves that are straight lines imply budget shares that
converge to constants as budgets grow and that leave the unit interval at low budgets. Household budget
surveys favor the Working-Leser form instead, in which budget shares are linear in the logarithm of the
budget (\cref{sec:05-engel}). Under that form the marginal propensity to spend changes with the budget,
as panel (b) of \cref{fig:05-gorman-piglog} shows, so a transfer between households moves aggregate
demand. We therefore want to relax linear Engel curves while keeping as much of the representative
consumer as possible.

Nonnegativity of demand breaks linear aggregation even when every household has quasi-homothetic
preferences \citep[\S 6.1, pp.~150--153]{DeatonMuellbauer1980b}. Fix prices and write a household's
Engel curve for one good as $q_h = d + cI_h$ with $c > 0$. If $d < 0$, the line implies negative demand
below the budget $I^* = -\frac{d}{c}$. Actual demand is zero there, so the true Engel curve is the
kinked line $q_h = \max\{0, c(I_h - I^*)\}$. A transfer from a household above the kink to one below it
lowers market demand, because the two have marginal propensities $c$ and $0$. Market demand then depends
on how many households sit below the kink. To see the shape, let budgets be uniform on $[\bar I -
\frac{r}{2}, \bar I + \frac{r}{2}]$, with $I^*$ inside this interval. Mean demand is
\begin{equation*}
  \bar q = \frac{1}{r}\int_{I^*}^{\bar I + \frac{r}{2}} c\,(I - I^*)\,\dd I
         = \frac{c}{2r}\Bigl(\bar I + \frac{r}{2} - I^*\Bigr)^2 ,
\end{equation*}
which is quadratic in mean expenditure. Its slope $\frac{c}{r}(\bar I + \frac{r}{2} - I^*)$ rises with
$\bar I$ as households cross the kink, and it reaches $c$ once all of them buy. This is the quadratic
case of \citet[Exercise~6.5]{DeatonMuellbauer1980b}, written in our notation.

The argument matters most for narrowly defined goods. Many consumers buy none of a particular brand, so
much of the market response to a higher budget comes from new buyers entering the market. Deaton and
Muellbauer call this the ``extensive'' margin, as opposed to the ``intensive'' margin of buyers who buy
more (p.~153). It would be a coincidence if the household Engel curves had just the offsetting
nonlinearity that made market demand linear. The discrete-choice models of \cref{ch:discrete,ch:logit}
put this extensive margin at the center, since each consumer there buys one product or none.

\begin{figure}[htbp]
\centering
\begin{subfigure}[t]{0.48\textwidth}
\centering
\begin{tikzpicture}
\begin{axis}[width=0.96\linewidth, height=0.84\linewidth,
  xmin=0, xmax=10.5, ymin=0, ymax=4.8,
  xlabel={total expenditure $I_h$}, ylabel={quantity of good 1, $q_{1h}$},
  xlabel style={font=\small}, ylabel style={font=\small}, tick label style={font=\scriptsize},
  axis lines=left, legend pos=north west, legend cell align=left,
  legend style={font=\scriptsize, draw=none, fill=none}]
  \addplot[thick, domain=1.5:10, samples=2] {0.55 + 0.3*x};
  \addplot[thick, dashed, domain=2.5:10, samples=2] {1.25 + 0.3*x};
  \addplot[only marks, mark=*, mark size=1.4pt, forget plot] coordinates {(1.5,1.0) (2.5,2.0)};
  \legend{household A, household B}
  \node[font=\scriptsize, align=left, anchor=north west] at (axis cs:3.6,1.3)
     {common slope $\frac{\partial q_{1h}}{\partial I_h} = \frac{m_1}{p_1}$};
\end{axis}
\end{tikzpicture}
\caption{Gorman polar form (linear expenditure system).}
\end{subfigure}\hfill
\begin{subfigure}[t]{0.48\textwidth}
\centering
\begin{tikzpicture}
\begin{axis}[width=0.96\linewidth, height=0.84\linewidth,
  xmin=0, xmax=8.5, ymin=0, ymax=2.8,
  xlabel={total expenditure $I_h$}, ylabel={quantity of good 1, $q_{1h}$},
  xlabel style={font=\small}, ylabel style={font=\small}, tick label style={font=\scriptsize},
  axis lines=left, legend pos=north west, legend cell align=left,
  legend style={font=\scriptsize, draw=none, fill=none}]
  \addplot[thick, domain=0.02:8, samples=120] {x*(0.5 - 0.10*ln(x))};
  \addplot[thin, dashed, domain=0:3.5, samples=2] {0.1 + 0.4*x};
  \addplot[thin, dashed, domain=1:8, samples=2, forget plot] {0.3 + 0.290139*x};
  \addplot[only marks, mark=*, mark size=1.4pt, forget plot] coordinates {(1,0.5) (3,1.170417)};
  \legend{AIDS Engel curve, tangents}
  \node[font=\scriptsize, anchor=north west] at (axis cs:1.05,0.45) {slope 0.40};
  \node[font=\scriptsize, anchor=north west] at (axis cs:3.05,1.12) {slope 0.29};
\end{axis}
\end{tikzpicture}
\caption{AIDS at unit prices (running example).}
\end{subfigure}
\caption{Engel curves under exact and almost exact aggregation. Panel (a): two households with the
linear expenditure system of \cref{ex:05-les}, $p = (1,1)$, $m_1 = 0.3$, subsistence quantities
$(\underline{q}_{1h}, \underline{q}_{2h}) = (1, 0.5)$ for household A (solid) and $(2, 0.5)$ for household
B (dashed). The lines start at the subsistence budgets (dots) and are parallel. Panel (b): the Engel curve
of good 1 in the running example of \cref{ex:05-running} at $p = (1,1,1)$, $q_1 = I(0.5 - 0.10\ln I)$. The
dashed tangents at the budgets $I = 1$ and $I = 3$ of the two households in \cref{ex:05-aggregation} have
slopes $0.4 - 0.10\ln I$, equal to 0.40 and 0.29. These are the households' marginal propensities to spend
on good 1. A transfer from the richer to the poorer household therefore raises aggregate demand for good
1. \Cref{fig:05-engel} draws the same curve together with a luxury.}
\label{fig:05-gorman-piglog}
\end{figure}

\subsection{Muellbauer's generalization: PIGL and PIGLOG}\label{sec:05-pigl}

\citet{Muellbauer1975,Muellbauer1976} weakens the requirement. The aggregate budget share of good $i$ is
aggregate spending on $i$ divided by aggregate spending. It equals the expenditure-weighted mean of
household shares:
\begin{equation}\label{eq:05-wbar}
  \bar w_i \equiv \frac{\sum_h p_i q_{ih}}{\sum_h I_h} = \sum_h \frac{I_h}{H\bar I}\, w_{ih}(p, I_h).
\end{equation}
Muellbauer's representative consumer exists if there is a scalar $I_0$, a function of prices and of the
distribution $(I_1, \dots, I_H)$, such that
\begin{equation}\label{eq:05-muellbauer}
  \bar w_i = w_i\bigl(p, I_0(p, I_1, \dots, I_H)\bigr) \quad \text{for every good } i.
\end{equation}
Here $w_i(p, \cdot)$ is the share function of a reference household
\citep[eq.~A1]{DeatonMuellbauer1980a}. We call $I_0$ the \emph{representative budget level}. The Gorman
case has $I_0 = \bar I$. Allowing $I_0$ to differ from the mean lets features of the distribution other
than its mean enter aggregate demand.

Muellbauer shows that \eqref{eq:05-muellbauer} holds if and only if household shares have the
\emph{generalized linear} form
\begin{equation}\label{eq:05-gl}
  w_{ih}(p, I_h) = v_h(p, I_h)\, A_i(p) + B_i(p) + \frac{C_{ih}(p)}{I_h},
\end{equation}
with $\sum_i A_i = \sum_i C_{ih} = \sum_h C_{ih} = 0$ and $\sum_i B_i = 1$
\citep[eq.~A2]{DeatonMuellbauer1980a}. The household-specific term $C_{ih}$ is an amount of spending, so it
enters the share divided by $I_h$. In the weighted mean \eqref{eq:05-wbar} it contributes $\frac{1}{H\bar
I}\sum_h C_{ih} = 0$.\footnote{\citet[eq.~A2]{DeatonMuellbauer1980a} print the term as $C_{ih}(p)$ with
$\sum_h C_{ih} = 0$. Read literally as a share term, it would leave $\sum_h \frac{I_h}{H\bar I}\,
C_{ih}$ in \eqref{eq:05-wbar}, which vanishes only if $\sum_h I_h C_{ih} = 0$. We write the term over
$I_h$, which is the reading under which it cancels.} We drop the household-specific terms $C_{ih}$ from
here on.

The preferences behind \eqref{eq:05-gl} have a simple structure \citep[\S 6.2, eqs.~2.3--2.6,
pp.~154--155]{DeatonMuellbauer1980b}. Household expenditure functions take the form $e_h(p, u_h) =
\theta_h[u_h, a(p), b(p)] + \phi_h(p)$. Here $a(p)$ and $b(p)$ are linearly homogeneous price indices
common to all households, and each $\theta_h$ is linearly homogeneous in $(a, b)$. The terms $\phi_h$ must
sum to zero over households, so the representative expenditure function is $e(p, u_0) = \theta[u_0,
a(p), b(p)]$. We can think of $a$ and $b$ as the prices of two intermediate goods. Shephard's lemma in
share form and the chain rule give
\begin{equation*}
  \bar w_i = \frac{\partial\ln\theta}{\partial\ln a}\,\frac{\partial\ln a}{\partial\ln p_i}
           + \frac{\partial\ln\theta}{\partial\ln b}\,\frac{\partial\ln b}{\partial\ln p_i} .
\end{equation*}
Euler's theorem for the degree-one function $\theta$ gives $\frac{\partial\ln\theta}{\partial\ln a} +
\frac{\partial\ln\theta}{\partial\ln b} = 1$. Writing $\varrho(I_0, p) = \frac{\partial\ln\theta}{\partial
\ln b}$ after substituting for $u_0$,
\begin{equation}\label{eq:05-gl-shares}
  \bar w_i = \bigl[1 - \varrho(I_0, p)\bigr]\frac{\partial\ln a}{\partial\ln p_i}
           + \varrho(I_0, p)\,\frac{\partial\ln b}{\partial\ln p_i} .
\end{equation}
Every aggregate share is a weighted average of the shares of the two indices, with one weight common to
all goods. Matching terms with \eqref{eq:05-gl} gives $B_i = \frac{\partial\ln a}{\partial\ln p_i}$, $A_i
= \frac{\partial\ln b}{\partial\ln p_i} - \frac{\partial\ln a}{\partial\ln p_i}$ and $v(p, I_0) =
\varrho(I_0, p)$. Hence
at constant prices each aggregate share is a linear function of every other share. With time-series data
in which relative prices barely move, this implication can be tested directly (p.~155).

The representative budget level is independent of prices if and only if
\begin{equation}\label{eq:05-pigl}
  v_h(p, I_h) = \frac{1 - \left(\frac{I_h}{\kappa_h}\right)^{-\psi}}{\psi}.
\end{equation}
Here $\psi$ is a constant and $\kappa_h > 0$ is a household-specific scale that does not depend on $I_h$
or $p$.\footnote{\citet[eq.~A3]{DeatonMuellbauer1980a} write the PIGL constant as $\alpha$. We use $\psi$
because $\alpha_i$ denotes the AIDS intercepts, and $\kappa_h$ for their $k_h$ because $k$ indexes goods.
The factor $\psi^{-1}$ appears in their typesetting as a superscript. We read it as a multiplicative
factor, which is the reading under which the limit $\psi \to 0$ gives their eq.~(A4).} Shares of this form
are called \emph{price-independent generalized linear} (PIGL). As $\psi \to 0$,
$\frac{1}{\psi}\bigl[1 - \exp\bigl(-\psi\ln\frac{I_h}{\kappa_h}\bigr)\bigr] \to \ln\frac{I_h}{\kappa_h}$, so
\begin{equation}\label{eq:05-piglog-v}
  v_h(p, I_h) = \ln \frac{I_h}{\kappa_h}.
\end{equation}
This limit is the \emph{price-independent generalized logarithmic} (PIGLOG) case.

\casehead{The representative budget level under PIGL}
With $C_{ih} = 0$, substituting \eqref{eq:05-gl} into \eqref{eq:05-wbar} gives
\begin{equation*}
  \bar w_i = A_i(p) \sum_h \frac{I_h}{H\bar I}\, v_h(p, I_h) + B_i(p).
\end{equation*}
A reference household with $\kappa = 1$ reproduces $\bar w_i$ if and only if $v(p, I_0) = \sum_h
\frac{I_h}{H\bar I} v_h(p, I_h)$. Under \eqref{eq:05-pigl} this condition reads
\begin{equation*}
  \frac{1 - I_0^{-\psi}}{\psi} = \sum_h \frac{I_h}{H\bar I}\,
  \frac{1 - \left(\frac{I_h}{\kappa_h}\right)^{-\psi}}{\psi} .
\end{equation*}
The weights $\frac{I_h}{H\bar I}$ sum to one, so the constant terms cancel and
\begin{equation}\label{eq:05-I0-pigl}
  I_0 = \left[\sum_h \frac{I_h}{H\bar I}\left(\frac{I_h}{\kappa_h}\right)^{-\psi}\right]^{-\frac{1}{\psi}} .
\end{equation}
The right-hand side involves budgets only, never prices. This is equation (A5) of
\citet{DeatonMuellbauer1980a}.\footnote{The version printed in their appendix omits the weight $I_h$
inside the sum. With the weight, $I_0$ equals the common budget when all budgets are equal and
$\kappa_h = 1$, as it must. Their book prints the formula with the weight. For $\kappa_h = 1$ it reads
$I_0 = \bigl[\frac{\sum_h I_h^{1-\psi}}{\sum_h I_h}\bigr]^{-\frac{1}{\psi}}$ \citep[\S 6.2, eq.~2.13,
p.~156]{DeatonMuellbauer1980b}, which is \eqref{eq:05-I0-pigl}.} Letting $\psi \to 0$ and expanding
$(\frac{I_h}{\kappa_h})^{-\psi} \approx 1 - \psi\ln\frac{I_h}{\kappa_h}$ gives the PIGLOG representative
budget,
\begin{equation}\label{eq:05-I0-piglog}
  \ln I_0 = \sum_h \frac{I_h}{H\bar I}\ln\frac{I_h}{\kappa_h}.
\end{equation}

The curvature parameter $\psi$ fixes how the distribution of budgets enters $I_0$. Take $\kappa_h = 1$.
Then \eqref{eq:05-I0-pigl} says that $I_0$ is a power mean of order $-\psi$ of the budgets, with
expenditure weights $\frac{I_h}{H\bar I}$. Power means rise with their order, so $I_0$ falls as $\psi$
rises. At $\psi = 1$, the case of linear Engel curves, the weighted mean of order $-1$ is
\begin{equation*}
  I_0 = \Bigl[\sum_h \frac{I_h}{H\bar I}\cdot\frac{1}{I_h}\Bigr]^{-1} = \Bigl[\frac{H}{H\bar I}\Bigr]^{-1}
      = \bar I ,
\end{equation*}
as in the Gorman case. Hence, when budgets are unequal, $I_0 > \bar I$ exactly when $\psi < 1$, the
empirically relevant case. Two
values of $\psi$ give closed forms \citep[\S 6.2, eqs.~2.14--2.15, p.~157]{DeatonMuellbauer1980b}. At
$\psi = -1$ the shares \eqref{eq:05-pigl-shares} are linear in the budget, so spending on each good is
quadratic in it. Then
\begin{equation*}
  I_0 = \frac{\sum_h I_h^2}{\sum_h I_h} = \frac{H(\sigma^2 + \bar I^2)}{H\bar I}
      = \bar I\Bigl(1 + \frac{\sigma^2}{\bar I^2}\Bigr),
\end{equation*}
where $\sigma^2 = \frac{1}{H}\sum_h (I_h - \bar I)^2$ is the variance of budgets. The representative budget
exceeds the mean by a factor of one plus the squared coefficient of variation. In the PIGLOG limit,
\cref{sec:05-almost-exact} shows that $\frac{I_0}{\bar I} = \frac{H}{Z}$, where $\ln Z$ is Theil's entropy
measure.
Deaton and Muellbauer's book writes this ratio as $\kappa_0 = \frac{I_0}{\bar I}$. It is the inverse of the
index $\kappa = \frac{Z}{H}$ of their article and of \cref{prop:05-almost-exact}.

\casehead{Cost functions and budget shares}
If household behavior comes from preferences, the PIGL cost function takes the form
\begin{equation}\label{eq:05-pigl-cost}
  \left[\frac{e(p, u_h)}{\kappa_h}\right]^{\psi} = (1 - u_h)\, a(p)^{\psi} + u_h\, b(p)^{\psi},
\end{equation}
and as $\psi \to 0$ it becomes the PIGLOG form
\begin{equation}\label{eq:05-piglog-cost}
  \ln\frac{e(p, u_h)}{\kappa_h} = (1 - u_h) \ln a(p) + u_h \ln b(p)
\end{equation}
\citep[eqs.~A6--A7]{DeatonMuellbauer1980a}. Here $a(p)$ and $b(p)$ are linearly homogeneous and concave,
and $0 \le u_h \le 1$ except as discussed below. With $\kappa_h = 1$, \eqref{eq:05-piglog-cost} is the
PIGLOG family of \cref{sec:05-piglog}.

The budget shares follow from Shephard's lemma in share form, $w_i = \frac{\partial \ln e}{\partial \ln
p_i}$. Take $\kappa_h = 1$, drop the household subscript, and write $a_i = \frac{\partial \ln a}{\partial
\ln p_i}$ and $b_i = \frac{\partial \ln b}{\partial \ln p_i}$. Differentiating \eqref{eq:05-pigl-cost}
with respect to $\ln p_i$, and using $\frac{\partial e^{\psi}}{\partial \ln p_i} = \psi e^{\psi}
\frac{\partial \ln e}{\partial \ln p_i}$, gives
\begin{equation*}
  \psi e^{\psi}\frac{\partial \ln e}{\partial \ln p_i} = (1 - u)\,\psi a^{\psi} a_i + u\, \psi b^{\psi} b_i .
\end{equation*}
At the optimum $e = I$. Dividing by $\psi I^{\psi}$,
\begin{equation}\label{eq:05-pigl-shares}
  w_i = (1 - u)\left(\frac{a}{I}\right)^{\psi} a_i + u\left(\frac{b}{I}\right)^{\psi} b_i,
  \qquad u = \frac{I^{\psi} - a^{\psi}}{b^{\psi} - a^{\psi}} .
\end{equation}
In the PIGLOG limit,
\begin{equation}\label{eq:05-piglog-shares}
  w_i = (1 - u)\, a_i + u\, b_i, \qquad u = \frac{\ln I - \ln a}{\ln b - \ln a}
\end{equation}
\citep[eqs.~A8--A10]{DeatonMuellbauer1980a}.

These equations have a transparent reading. By \eqref{eq:05-pigl-cost}, $(1-u)a^{\psi} + ub^{\psi} =
I^{\psi}$, so the weights $(1-u)(\frac{a}{I})^{\psi}$ and $u(\frac{b}{I})^{\psi}$ sum to one. Budget
shares are weighted averages of $a_i$ and $b_i$. As utility rises from $0$ to $1$, cost rises from $a(p)$
to $b(p)$ whenever $b > a$, whatever the sign of $\psi$, and the shares move from $a_i$ to $b_i$.
\citet{DeatonMuellbauer1980a} call $a(p)$ ``poverty'' expenditure, with spending pattern $a_i$, and $b(p)$
``affluence'' expenditure, with pattern $b_i$. A good is a luxury if and only if $b_i > a_i$, because its
share then rises with $u$. Inferior goods are not excluded.

For the AIDS, \eqref{eq:05-aids-cost} gives $\ln a(p) = \ln\mathbf{P}(p)$ and $\ln b(p) = \ln a(p) +
\beta_0\prod_k p_k^{\beta_k}$. Because $\frac{\partial}{\partial \ln p_i}\prod_k p_k^{\beta_k} = \beta_i
\prod_k p_k^{\beta_k}$,
\begin{equation*}
  b_i - a_i = \beta_i\, \beta_0 \prod_k p_k^{\beta_k}.
\end{equation*}
With $\beta_0 > 0$, good $i$ is a luxury exactly when $\beta_i > 0$, as found in \cref{sec:05-aids}.

\casehead{Where the demands are valid}
Shares must lie in $[0, 1]$. Under PIGLOG each $w_i$ is linear in $u$, so the admissible values of $u$
form an interval. Because $a(p)$ and $b(p)$ are increasing and linearly homogeneous, $a_i \ge 0$, $b_i
\ge 0$ and $\sum_i a_i = \sum_i b_i = 1$. Hence $a_i, b_i \in [0,1]$, and every $u \in [0,1]$ is
admissible. Above $u = 1$, the share of a luxury reaches one at $u = \frac{1 - a_i}{b_i - a_i}$, and the
share of a necessity reaches zero at $u = \frac{a_i}{a_i - b_i}$. The upper end of the admissible
interval is the smallest of these values. Below $u = 0$ the same reasoning, with the roles of zero and
one exchanged, gives the lower end. \citet[p.~324]{DeatonMuellbauer1980a} state the corresponding limits
for the PIGL case.

In the running example at unit prices, $\ln a = \alpha_0 = 0$ and $\ln b - \ln a = \beta_0 = 1$, so $u =
\ln I$. The share vectors are $a = \alpha = (0.5, 0.3, 0.2)$ and $b = \alpha + \beta = (0.40, 0.34, 0.26)$.
The necessity, good 1, has its share reach zero at $u = \frac{0.5}{0.1} = 5$. Goods 2 and 3 reach a share
of one only at $u = 17.5$ and $u = 13.3$. Below zero, the share of good 3 reaches zero first, at $u =
-\frac{0.2}{0.06} = -3.33$. The demands are therefore valid for $\ln I \in [-3.33, 5]$, the range reported
for the running example in \cref{sec:05-aids-limits}.

\casehead{Concavity}
A sufficient condition for the PIGLOG cost function to be concave in $p$ is that $a$ and $b$ be concave
and $0 \le u \le 1$. To see this, write \eqref{eq:05-piglog-cost} as $e = a(p)^{1-u}\, b(p)^{u}$. The
weighted geometric mean $g(x, y) = x^{1-u}y^{u}$ is concave and nondecreasing on the positive orthant for
$u \in [0, 1]$. A concave nondecreasing function of concave functions is concave. The condition is not
necessary. \citet[p.~324]{DeatonMuellbauer1980a} show that the PIGL cost function is concave for
\emph{all} $I > 0$ and $p > 0$ only if $a_i = b_i$ for every $i$, in which case preferences are
homothetic. Any nonhomothetic PIGLOG system, the AIDS included, must therefore have its concavity checked
where it is used, as in \cref{prop:05-negativity}.

\casehead{Two special cases}
With $\psi = 1$, \eqref{eq:05-pigl-cost} reads $e_h = \kappa_h a(p) + \kappa_h u_h[b(p) - a(p)]$. Defining
$\tilde u_h = \kappa_h u_h$ gives $e_h = \kappa_h a(p) + [b(p) - a(p)]\tilde u_h$, a Gorman polar form with
$a_h = \kappa_h a$ and common slope function $b - a$. So the PIGL family contains every model with linear
Engel curves, such as the linear expenditure system \citep[p.~324]{DeatonMuellbauer1980a}. The second
case is less obvious.

\begin{remark}[When the indirect translog is PIGLOG]\label{rem:05-translog-piglog}
The indirect translog of \citet{ChristensenJorgensonLau1975} specifies indirect utility as
\begin{equation*}
  v = \tau_0 + \sum_i \tau_i \ln\frac{p_i}{I}
      + \frac{1}{2}\sum_i\sum_j \tau_{ij}\ln\frac{p_i}{I}\ln\frac{p_j}{I},
\end{equation*}
with the normalizations $\sum_i \tau_i = -1$ and $\tau_{ij} = \tau_{ji}$. Let $\tau_{Mi} = \sum_k
\tau_{ki}$ and $L = \ln I$. Expanding the products and collecting powers of $L$,
\begin{equation*}
  v = \tau_0 + \sum_i \tau_i \ln p_i + \frac{1}{2}\sum_i\sum_j \tau_{ij}\ln p_i \ln p_j
      + L\Bigl(1 - \sum_i \tau_{Mi}\ln p_i\Bigr) + \frac{1}{2}L^2 \sum_i \tau_{Mi} .
\end{equation*}
If $\sum_i \tau_{Mi} = 0$, utility is linear in $L$ and can be solved for the log cost function:
\begin{equation*}
  \ln e(p, v) = \frac{v - \tau_0 - \sum_i \tau_i \ln p_i - \frac{1}{2}\sum_i\sum_j\tau_{ij}\ln p_i\ln p_j}
                     {1 - \sum_i \tau_{Mi}\ln p_i} .
\end{equation*}
This is linear in $v$, with an intercept and a slope that depend on prices, which is the PIGLOG form
\eqref{eq:05-piglog-cost}. If $\sum_i\tau_{Mi} \neq 0$, utility is quadratic in $L$ and $\ln e$ is not
linear in $v$. Utility is ordinal, so this alone does not exclude a PIGLOG form for some increasing
transformation of $v$. Budget shares do not depend on how utility is normalized, and they settle the
question. Write $\tau_{MM} = \sum_i\tau_{Mi}$, $N_i(p) = \tau_i + \sum_j\tau_{ij}\ln p_j$ and $D(p) =
\sum_j\tau_{Mj}\ln p_j - 1$. Roy's identity and the symmetry of $\tau_{ij}$ give
\begin{equation*}
  w_i = \frac{N_i(p) - \tau_{Mi}L}{D(p) - \tau_{MM}L}
      = \frac{\tau_{Mi}}{\tau_{MM}} + \frac{c_i(p)}{D(p) - \tau_{MM}L},
  \qquad c_i(p) = N_i(p) - \frac{\tau_{Mi}}{\tau_{MM}}D(p),
\end{equation*}
where the second form requires $\tau_{MM} \neq 0$. By \eqref{eq:05-pigl-shares}, PIGL shares are affine
in $I^{-\psi}$, and PIGLOG shares are affine in $L$. The last term has neither form on an interval of
budgets unless $c_i(p) = 0$. Matching the constant and the coefficients of $\ln p_j$ in $c_i(p) = 0$ gives
$\tau_{Mi} = -\tau_i\tau_{MM}$ and $\tau_{ij} = \tau_i\tau_j\tau_{MM}$, so that $w_i = -\tau_i$ at all
prices and budgets. These are Cobb-Douglas preferences. With $t = L + \sum_i\tau_i\ln p_i$, utility is
$v = \tau_0 + t + \frac{1}{2}\tau_{MM}t^2$, which represents them only where $1 + \tau_{MM}t > 0$.
Apart from this homothetic case, the indirect translog is PIGLOG exactly when $\tau_{MM} = 0$. For
$\tau_{MM} \neq 0$ the display still has the generalized linear form \eqref{eq:05-gl}, with
$v_h(p, I_h) = \frac{1}{D(p) - \tau_{MM}\ln I_h}$, $A_i = c_i(p)$ and $B_i = \frac{\tau_{Mi}}{\tau_{MM}}$.
Because $\sum_i N_i(p) = D(p)$, these satisfy $\sum_i c_i(p) = 0$ and $\sum_i B_i = 1$. A
representative budget level in the sense of \eqref{eq:05-muellbauer} then exists. The shares are
neither PIGL nor PIGLOG, so by \eqref{eq:05-pigl} that level depends on prices.
\citet[p.~325, eqs.~A11--A12]{DeatonMuellbauer1980a} call $\tau_{MM} = 0$ the only case in which the
indirect translog allows consistent aggregation. Their claim is right for a representative budget level
that does not depend on prices, apart from the homothetic exception, which they do not note. They add
that no such result holds for any interesting special case of the direct translog.
\end{remark}

\casehead{Evidence on the curvature parameter}
\citet[\S 6.2, p.~158]{DeatonMuellbauer1980b} fit the PIGL demand functions to postwar British data for
1954--1974 with eight groups of goods, the same period and number of groups as in
\cref{sec:05-dm-application}. They take $a(p) = \sum_k
a_kp_k$, the cost of a fixed subsistence bundle, and $b(p) = \beta_0\prod_k p_k^{\beta_k}$. The estimate of
$\psi$ is very close to $-1$, which implies nearly quadratic Engel curves over time. The hypothesis $\psi =
1$, which gives linear Engel curves, is decisively rejected. The book states these results in words and
reports no numerical estimate. A full use of the PIGL class would also measure $I_0$ from the
distribution of budgets. Few countries have data on the distribution of total expenditure of sufficient
quality for this, and the authors' own indirect attempts had no notable success. Applied work therefore
replaces $I_0$ by $\bar I$. The error is small if the distribution is stable, so that $\frac{I_0}{\bar I}$
is roughly constant, which they judge possible for postwar Britain (p.~158).
\Cref{sec:05-almost-exact} shows what the approximation costs when the distribution moves.

\subsection{Almost exact aggregation of the AIDS}\label{sec:05-almost-exact}

The AIDS is a PIGLOG system, so it aggregates through a single representative budget.
\citet[Section~I.A]{DeatonMuellbauer1980a} allow households to differ through a scale $\kappa_h$. They
interpret it as ``a sophisticated measure of household size'' that could reflect age composition, other
characteristics and economies of household size. Household $h$ deflates its budget by $\kappa_h$ to a
``needs corrected'' per capita level.

\begin{proposition}[Almost exact aggregation]\label{prop:05-almost-exact}
Suppose every household has the AIDS preferences of \eqref{eq:05-aids-cost} with common parameters, except
for a scale $\kappa_h > 0$ that deflates its budget, so that
\begin{equation}\label{eq:05-aids-household}
  w_{ih} = \alpha_i + \sum_j \gamma_{ij}\ln p_j + \beta_i \ln\frac{I_h}{\kappa_h \mathbf{P}(p)} .
\end{equation}
Define the aggregate index $\kappa$ by
\begin{equation}\label{eq:05-kappa}
  \ln\frac{\bar I}{\kappa} = \sum_h \frac{I_h}{H\bar I}\ln\frac{I_h}{\kappa_h}.
\end{equation}
Then the aggregate budget shares are
\begin{equation}\label{eq:05-aids-aggregate}
  \bar w_i = \alpha_i + \sum_j \gamma_{ij}\ln p_j + \beta_i \ln\frac{\bar I}{\kappa\, \mathbf{P}(p)} ,
\end{equation}
the shares of a single AIDS household with budget $\frac{\bar I}{\kappa}$. When $\kappa_h = 1$ for all
households, $\frac{\bar I}{\kappa} = I_0$ with
\begin{equation}\label{eq:05-I0}
  \ln I_0 = \frac{\sum_h I_h \ln I_h}{\sum_h I_h} .
\end{equation}
\end{proposition}

\begin{proof}
Let $A_i(p) = \alpha_i + \sum_j\gamma_{ij}\ln p_j - \beta_i\ln\mathbf{P}(p)$, so that $w_{ih} = A_i(p) +
\beta_i\ln\frac{I_h}{\kappa_h}$. Substitute into \eqref{eq:05-wbar}:
\begin{align*}
  \bar w_i &= \sum_h \frac{I_h}{H\bar I}\Bigl[A_i(p) + \beta_i\ln\frac{I_h}{\kappa_h}\Bigr] \\
           &= A_i(p)\sum_h \frac{I_h}{H\bar I} + \beta_i \sum_h \frac{I_h}{H\bar I}\ln\frac{I_h}{\kappa_h} \\
           &= A_i(p) + \beta_i\ln\frac{\bar I}{\kappa} .
\end{align*}
The last line uses the fact that the weights sum to one and the definition \eqref{eq:05-kappa}.
Substituting $A_i(p)$ back gives \eqref{eq:05-aids-aggregate}. With $\kappa_h = 1$, \eqref{eq:05-kappa}
reads $\ln\frac{\bar I}{\kappa} = \sum_h \frac{I_h}{H\bar I}\ln I_h$, which is \eqref{eq:05-I0}.
\end{proof}

The same argument works for any PIGLOG system. From \eqref{eq:05-piglog-shares}, $u_h = \frac{\ln I_h -
\ln a}{\ln b - \ln a}$, so
\begin{equation*}
  w_{ih} = a_i + (b_i - a_i)u_h
  = \Bigl[a_i - \frac{b_i - a_i}{\ln b - \ln a}\ln a\Bigr] + \frac{b_i - a_i}{\ln b - \ln a}\ln I_h .
\end{equation*}
Household shares are linear in $\ln I_h$ with price-dependent coefficients $A_i(p)$ and $B_i(p)$. For the
AIDS, $b_i - a_i = \beta_i(\ln b - \ln a)$, so $B_i = \beta_i$ and $A_i = a_i - \beta_i\ln a = \alpha_i +
\sum_j\gamma_{ij}\ln p_j - \beta_i\ln\mathbf{P}$, as in the proof. Aggregation is ``almost'' exact
because aggregate demand depends on the distribution of budgets, but only through one scalar index.

\casehead{What $I_0$ measures}
Write each budget as $I_h = \bar I\cdot\frac{I_h}{\bar I}$. Then
\begin{align}
  \ln I_0 &= \sum_h \frac{I_h}{H\bar I}\Bigl[\ln\bar I + \ln\frac{I_h}{\bar I}\Bigr]
           = \ln\bar I + \mathcal{T}, \notag\\
  \mathcal{T} &= \frac{1}{H}\sum_h \frac{I_h}{\bar I}\ln\frac{I_h}{\bar I}. \label{eq:05-theil}
\end{align}
The quantity $\mathcal{T}$ is the \emph{Theil index} of inequality. The function $f(r) = r\ln r$ is
convex and the ratios $r_h = \frac{I_h}{\bar I}$ average one. Jensen's inequality therefore gives
$\mathcal{T} = \frac{1}{H}\sum_h f(r_h) \ge f(1) = 0$, with equality only when all budgets are equal. The
maximum, $\ln H$, is reached when one household does all the spending. Hence
\begin{equation*}
  I_0 = \bar I\, e^{\mathcal{T}} \ge \bar I .
\end{equation*}
The representative household is richer than the average household, because expenditure weighting gives
rich households more weight.

Deaton and Muellbauer express the same object through an entropy measure. With $\kappa_h = 1$ they write
$\kappa = \frac{Z}{H}$, where $\ln Z = -\sum_h \frac{I_h}{H\bar I}\ln\frac{I_h}{H\bar I}$ is Theil's
entropy measure of equality. Since $\ln\frac{I_h}{H\bar I} = \ln\frac{I_h}{\bar I} - \ln H$,
\begin{equation*}
  \ln Z = \ln H - \mathcal{T}, \qquad \ln\kappa = \ln Z - \ln H = -\mathcal{T}, \qquad
  \frac{\bar I}{\kappa} = \bar I e^{\mathcal{T}} = I_0 .
\end{equation*}
Hence $\frac{I_0}{\bar I} = \frac{H}{Z}$, which is the ratio $\kappa_0$ of \citet[\S 6.2, eq.~2.15,
p.~157]{DeatonMuellbauer1980b}. The entropy $Z$ reaches its maximum $H$ under perfect equality. In their words, ``as inequality
increases, $k = \frac{Z}{H}$ decreases and the representative budget level increases''
\citep[p.~315]{DeatonMuellbauer1980a}. When the scales $\kappa_h$ differ across households, $\kappa$
reflects the demographic structure as well as the distribution of budgets.

\casehead{Consequences for estimation}
In a time series of aggregate shares, \eqref{eq:05-aids-aggregate} reads
\begin{equation}\label{eq:05-aids-agg-ts}
  \bar w_{it} = \alpha_i + \sum_j\gamma_{ij}\ln p_{jt} + \beta_i\ln\frac{\bar I_t}{\mathbf{P}_t}
                - \beta_i\ln\kappa_t .
\end{equation}
Researchers usually observe mean expenditure $\bar I_t$ but not $\kappa_t$, which requires the joint
distribution of budgets and household characteristics. If $\kappa_t$ is constant, or varies independently
of $\bar I_t$ and $p_t$, we can replace $\ln\kappa_t$ by a constant $\ln\kappa^*$, its value or sample
mean. The intercepts then become $\alpha_i^* = \alpha_i - \beta_i\ln\kappa^*$
\citep[p.~315]{DeatonMuellbauer1980a}. The deviation $-\beta_i(\ln\kappa_t - \ln\kappa^*)$ joins the
error term and is uncorrelated with the regressors.

If $\kappa_t$ moves with the regressors, omitting it biases the estimates. Take the case $\kappa_h = 1$,
so that $-\ln\kappa_t = \mathcal{T}_t$, and hold prices fixed for simplicity. Let the linear projection of
inequality on log mean expenditure be $\mathcal{T}_t = c + \delta_{\mathcal{T}}\ln\bar I_t + r_t$, with
$r_t$ uncorrelated with $\ln\bar I_t$. Substituting into \eqref{eq:05-aids-agg-ts} and adding an error
term $\xi_{it}$,
\begin{equation*}
  \bar w_{it} = \bigl(\alpha_i - \beta_i\ln\mathbf{P} + \beta_i c\bigr)
               + \beta_i(1 + \delta_{\mathcal{T}})\ln\bar I_t + \beta_i r_t + \xi_{it}.
\end{equation*}
A regression of $\bar w_{it}$ on $\ln\bar I_t$ estimates $\beta_i(1 + \delta_{\mathcal{T}})$. If
inequality rises with mean expenditure ($\delta_{\mathcal{T}} > 0$), luxuries look more luxurious and
necessities more necessary than they are. With several regressors the same logic applies to the partial
correlation of $\mathcal{T}_t$ with each of them. Deaton and Muellbauer list a non-constant $\kappa$
correlated with $\bar I$ and $p$ as one possible reason for the rejection of homogeneity in their data
(\cref{sec:05-dm-application}).

The index $\kappa$ is constant when the distribution of budgets and household characteristics changes
only through equiproportional changes in all budgets. Deaton and Muellbauer point out that in this case
even $\alpha_{ih}$ and $\gamma_{ijh}$ may differ across households. The reason is that the weights
$\frac{I_h}{H\bar I}$ are then constant over time, and each household's $\ln\mathbf{P}_h$ is linear in its
own parameters. With a common $\beta_i$, the weighted sum of the household equations is an AIDS whose
$\alpha_i$ and $\gamma_{ij}$ are expenditure-weighted averages of the household parameters.

\begin{example}[Two households]\label{ex:05-aggregation}
Take the running example of \cref{ex:05-running} at $p = (1,1,1)$, where $\ln\mathbf{P} = 0$ and $w_{ih} =
\alpha_i + \beta_i\ln I_h$. Two households have budgets $I_1 = 1$ and $I_2 = 3$, so $\bar I = 2$. The
expenditure weights are $\frac{1}{4}$ and $\frac{3}{4}$, and
\begin{equation*}
  \ln I_0 = \frac{1}{4}\ln 1 + \frac{3}{4}\ln 3 = 0.8240, \qquad I_0 = 2.280 .
\end{equation*}
The Theil index is $\mathcal{T} = \ln I_0 - \ln\bar I = 0.8240 - 0.6931 = 0.1308$. Equivalently,
$\kappa = e^{-\mathcal{T}} = 0.877$ and $Z = H\kappa = 1.755$. \Cref{tab:05-agg-example} reports the
household shares, the aggregate shares computed from \eqref{eq:05-wbar}, and the shares of a single
consumer at $I_0$ and at $\bar I$. The consumer with budget $I_0$ reproduces the aggregate shares
exactly. The consumer with the mean budget overstates the share of the necessity, good 1, by 1.3
percentage points and understates the share of the luxury, good 3, by 0.8 points.
\end{example}

\begin{table}[htbp]
\centering
\caption{Aggregation in the running example at unit prices}
\label{tab:05-agg-example}
\begin{tabular}{lccccc}
\toprule
 & Budget & $\ln$ budget & $w_1$ & $w_2$ & $w_3$ \\
\midrule
Household 1 & 1.000 & 0.0000 & 0.5000 & 0.3000 & 0.2000 \\
Household 2 & 3.000 & 1.0986 & 0.3901 & 0.3439 & 0.2659 \\
Aggregate, \eqref{eq:05-wbar} & 4.000 & & 0.4176 & 0.3330 & 0.2494 \\
Consumer at $I_0$ & 2.280 & 0.8240 & 0.4176 & 0.3330 & 0.2494 \\
Consumer at $\bar I$ & 2.000 & 0.6931 & 0.4307 & 0.3277 & 0.2416 \\
\bottomrule
\end{tabular}

\smallskip
\begin{minipage}{0.82\textwidth}
\footnotesize Notes: parameters of \cref{ex:05-running}, $w_{ih} = \alpha_i + \beta_i\ln I_h$ at $p =
(1,1,1)$. The aggregate row gives total expenditure and expenditure-weighted shares. Computed by
\texttt{figures/ch05/src/aggregation\_example.py}.
\end{minipage}
\end{table}

\begin{remark}[Welfare with a representative household]\label{rem:05-agg-welfare}
Under the Gorman polar form the representative consumer is normative: its compensating variation is the
sum of household compensating variations (\cref{prop:05-gorman}). Under PIGLOG the household with budget
$I_0$ is a positive construct, and its expenditure function does not aggregate the household expenditure
functions. How large is the error if we use it anyway? For the AIDS, the initial utility of household $h$
is $u_h^0 = \frac{\ln I_h - \ln\mathbf{P}(p^0)}{B(p^0)}$ by \eqref{eq:05-aids-v}, with $B(p) =
\beta_0\prod_k p_k^{\beta_k}$ as in \eqref{eq:05-aids-B}. Its cost of reaching $u_h^0$ at new prices $p^1$
satisfies
\begin{equation*}
  \ln e(p^1, u_h^0) = \ln\mathbf{P}(p^1) + g\bigl[\ln I_h - \ln\mathbf{P}(p^0)\bigr],
  \qquad g = \frac{B(p^1)}{B(p^0)} ,
\end{equation*}
so $e(p^1, u_h^0) = \mathbf{P}(p^1)\mathbf{P}(p^0)^{-g} I_h^{g}$. Summing over households gives
$\mathbf{P}(p^1)\mathbf{P}(p^0)^{-g}\sum_h I_h^g$. A natural representative calculation computes the
compensating variation of the household with budget $I_0$ as a fraction of its budget and applies that
fraction to aggregate expenditure $H\bar I$. This gives $\mathbf{P}(p^1)\mathbf{P}(p^0)^{-g} H\bar I\,
I_0^{g-1}$ for the aggregate cost at new prices. Since the exponential function is convex, Jensen's
inequality gives
\begin{equation*}
  \sum_h I_h^{g} = H\bar I\sum_h \frac{I_h}{H\bar I}\, e^{(g-1)\ln I_h}
  \ge H\bar I\, e^{(g-1)\sum_h \frac{I_h}{H\bar I}\ln I_h} = H\bar I\, I_0^{g-1}.
\end{equation*}
The representative calculation therefore understates the sum of household compensating variations. The
two coincide when $g = 1$ or when all budgets are equal, and their first-order terms in $g - 1$ agree, so
the gap is of second order. In \cref{ex:05-aggregation}, a 10 percent rise in $p_3$ gives $g = 1.1^{0.06}
= 1.0057$. Household compensating variations are 0.0184 (1.84 percent of the budget) and 0.0745 (2.48
percent). Their sum is 0.09296, against 0.09294 from the representative calculation. For a 50 percent
rise the figures are 0.3621 and 0.3618. The aggregate is accurate, but it hides who bears the loss. The
richer household loses a larger fraction of its budget because good 3 is a luxury. The welfare numbers
of \citet{ChaudhuriGoldbergJia2006}, studied in \cref{sec:05-cgj}, are representative-household numbers
of this kind. The authors state that their framework is not suited to asking how different groups in the
population are affected (their Section 3.6).\footnote{Page, table, equation and footnote numbers for
\citet{ChaudhuriGoldbergJia2006} refer throughout this chapter to the revised NBER working paper (Working Paper
10159, revised July 2006). Page numbers are those printed in it. Where we cite the published article instead, the
text says so.} They add that this
question would need micro data on purchases (their fn.~30). The same issue arises for price indices.
Under PIGLOG the exact cost-of-living index differs across households whenever the prices of luxuries
and necessities move apart. \Cref{sec:05-welfare-cv} derives the household index and reports British
estimates in which inflation in 1975--76 ranged from 16.39 percent for a poor couple to 14.05 percent for
a rich one (\cref{tab:05-dm-table71}).
\end{remark}

\section{Separability and multi-stage budgeting}\label{sec:05-separability}

A product-space system for $J$ goods has of the order of $J^2$ price coefficients
(\cref{sec:05-aids-limits}). With a hundred brands the AIDS has 4,950 free price coefficients even after
symmetry and homogeneity, far more than any data set can pin down. The classical remedy is to arrange
goods in a tree. The consumer first allocates total expenditure across broad groups and then allocates
each group's budget across the goods in the group. This section states the preferences under which such
\emph{multi-stage budgeting} is consistent with utility maximization, what they imply for substitution
between groups, and how elasticities at different levels of the tree combine.

Throughout, the goods are partitioned into groups $G = 1, \dots, \mathcal{G}$. We write $q_G$ and $p_G$
for the quantities and prices of the goods in group $G$, $X_G = p_G^{\T} q_G$ for group expenditure, and
$w_{i|G} = \frac{p_i q_i}{X_G}$ for the \emph{within-group share} of good $i \in G$. A second group is
denoted $G'$.

\subsection{Weak separability and the second stage}\label{sec:05-second-stage}

\begin{definition}[Weak separability]\label{def:05-weak-sep}
Preferences are \emph{weakly separable} in the partition $\{G\}$ if the utility function can be written as
\begin{equation}\label{eq:05-weak-sep}
  U(q) = F\bigl(U_1(q_1), U_2(q_2), \dots, U_{\mathcal{G}}(q_{\mathcal{G}})\bigr).
\end{equation}
Here $F$ is strictly increasing in each argument, and each \emph{subutility function} $U_G$ is
continuous, increasing and strictly quasi-concave. Preferences are \emph{additively separable} across groups if, in
addition, $F$ can be chosen as an increasing function of $\sum_G U_G(q_G)$.
\end{definition}

Weak separability says that the marginal rate of substitution between two goods in the same group does
not depend on the quantities of goods outside the group. Hence tastes within a group can be described
without reference to the rest of the budget.

\begin{proposition}[The second stage]\label{prop:05-second-stage}
Under weak separability, the Marshallian demand for a good $i \in G$ can be written as
\begin{equation}\label{eq:05-conditional-demand}
  q_i(p, I) = q_{i|G}\bigl(p_G, X_G(p, I)\bigr),
\end{equation}
where $X_G(p, I)$ is the optimal expenditure on group $G$ and $q_{i|G}(p_G, X_G)$ solves
$\max_{q_G} U_G(q_G)$ subject to $p_G^{\T} q_G = X_G$. The \emph{conditional demand} functions
$q_{i|G}$ depend only on the prices of the group and on group expenditure.
\end{proposition}

\begin{proof}
Let $q^*$ be the optimal bundle at $(p, I)$ and let $X_G^* = p_G^{\T} q_G^*$. Suppose $q_G^*$ did not
solve the group problem at budget $X_G^*$. Then some $q_G'$ with $p_G^{\T} q_G' \le X_G^*$ would give
$U_G(q_G') > U_G(q_G^*)$. Replacing $q_G^*$ by $q_G'$ in $q^*$ keeps total spending within $I$ and, because
$F$ is strictly increasing, raises $U$. This contradicts the optimality of $q^*$. So $q_G^*$ solves the
group problem, whose solution is unique by strict quasi-concavity of $U_G$. Hence $q_G^* =
q_G(p_G, X_G^*)$.
\end{proof}

The conditional demands are ordinary demand functions generated by the subutility $U_G$, with $X_G$ in
the role of income. They add up to $X_G$, are homogeneous of degree zero in $(p_G, X_G)$, and have a
symmetric and negative semidefinite conditional Slutsky matrix. Any flexible demand system, the AIDS in
particular, can therefore be fitted to within-group shares $w_{i|G}$ with group expenditure as the
budget variable. This is the bottom level of both multi-level systems studied in this chapter. It gives
the quinolone equations of \citet{ChaudhuriGoldbergJia2006} (\cref{sec:05-cgj-demand}) and the brand
equations of \citet{HausmanLeonardZona1994} and \citet{Hausman1997} (\cref{sec:05-hlz}). Hausman
stresses that the budget variable must be segment expenditure. Using overall expenditure instead ``is
inconsistent with the economic theory of multistage budgeting, and it can lead to decidedly inferior
econometric results'' \citep[p.~217]{Hausman1997}.

\subsection{The first stage}\label{sec:05-first-stage}

Weak separability alone does not deliver a simple first stage. Group expenditure $X_G(p, I)$ depends, in
general, on every individual price in every group. Multi-stage budgeting in the useful sense requires
that the allocation of total expenditure depend on group prices only through group price indices
$P_G(p_G)$, so that $X_G = X_G(P_1, \dots, P_{\mathcal{G}}, I)$. Two preference structures deliver
this. The second uses the following functional form.

\begin{definition}[Generalized Gorman polar form]\label{def:05-ggpf}
A sub-indirect utility function $V_G(p_G, X_G) = \max\{U_G(q_G) : p_G^{\T}q_G \le X_G\}$ has the
\emph{generalized Gorman polar form} (GGPF) if
\begin{equation}\label{eq:05-ggpf}
  V_G(p_G, X_G) = F_G\!\left(\frac{X_G}{\mathcal{P}_G(p_G)}\right) + \mathcal{A}_G(p_G).
\end{equation}
Here $F_G$ is increasing, $\mathcal{P}_G$ is homogeneous of degree one and $\mathcal{A}_G$ is homogeneous
of degree zero in $p_G$.
\end{definition}

\begin{proposition}[Two-stage budgeting, \citealp{Gorman1959}]\label{prop:05-first-stage}
Suppose preferences are weakly separable.
\begin{enumerate}[label=(\roman*)]
  \item If every subutility $U_G$ is homothetic, group expenditures depend on prices only through the
  group price indices $P_G(p_G) = e_G(p_G, 1)$, the unit costs of the subutilities.
  \item If preferences are additively separable across groups, and every sub-indirect utility function
  of the additive representation has the GGPF \eqref{eq:05-ggpf}, group expenditures depend on prices
  only through the indices $\mathcal{P}_G(p_G)$.
\end{enumerate}
\citet{Gorman1959} showed that, except when there are only two groups, these two structures are also
necessary for exact two-stage budgeting under weak separability \citep[\S 5.2,
p.~131]{DeatonMuellbauer1980b}. \citet{BlackorbyPrimontRussell1978} give a complete statement.
\end{proposition}

\begin{proof}
(i) Normalize each homothetic $U_G$ to be homogeneous of degree one. Its expenditure function is then
$e_G(p_G, u_G) = P_G(p_G)\, u_G$, with $P_G(p_G) = e_G(p_G, 1)$ linearly homogeneous and concave. By
\cref{prop:05-second-stage} the consumer chooses subutility levels to solve
\begin{equation*}
  \max_{u_1, \dots, u_{\mathcal{G}}} F(u_1, \dots, u_{\mathcal{G}})
  \quad\text{subject to}\quad \sum_G P_G(p_G)\, u_G \le I .
\end{equation*}
This is a standard consumer problem with ``quantities'' $u_G$ and ``prices'' $P_G$. Its solution $u_G^*$
depends only on $(P_1, \dots, P_{\mathcal{G}}, I)$, and so does $X_G = P_G u_G^*$.

(ii) Let $\sum_G U_G$ be the additive representation, and let $V_G$ be the sub-indirect utility of
$U_G$. The first stage maximizes $\sum_G V_G(p_G, X_G)$ subject to $\sum_G X_G = I$, after an increasing
transformation of the objective. At an interior optimum the marginal utilities
of group expenditure are equal:
\begin{equation*}
  \frac{\partial V_G}{\partial X_G} = \frac{1}{\mathcal{P}_G(p_G)}\,
  F_G'\!\left(\frac{X_G}{\mathcal{P}_G(p_G)}\right) = \lambda \quad \text{for every } G.
\end{equation*}
Each equation involves $p_G$ only through $\mathcal{P}_G$. Solving for $X_G$ in terms of $\lambda$ and
$\mathcal{P}_G$, and choosing $\lambda$ so that the $X_G$ sum to $I$, gives $X_G$ as a function of
$(\mathcal{P}_1, \dots, \mathcal{P}_{\mathcal{G}}, I)$. The term $\mathcal{A}_G$ does not affect the
allocation.
\end{proof}

Both routes are restrictive, in different places \citep[\S 5.2, pp.~130--131]{DeatonMuellbauer1980b}.
Route (i) restricts tastes within groups. With $e_G = P_G(p_G)\,u_G$, Shephard's lemma in share form gives
the within-group shares
\begin{equation*}
  w_{i|G} = \frac{\partial\ln e_G(p_G, u_G)}{\partial\ln p_i} = \frac{\partial\ln P_G(p_G)}{\partial\ln p_i},
\end{equation*}
which do not depend on $u_G$ and hence not on $X_G$. Every good in the group then has an elasticity of
one with respect to group expenditure. A group cannot mix necessities and luxuries, so a segment
containing both cheap and premium brands violates route (i). Route (ii) permits nonlinear within-group
Engel curves. Its price is additivity across groups, which Deaton and Muellbauer call ``both severe and,
in our view, unrealistic'' (p.~131). A reasonably general solution must therefore be approximate, or it
must abandon weak separability. The rest of this subsection takes the first path, and
\cref{rem:05-implicit-sep} sketches the second.

\casehead{The AIDS is a generalized Gorman polar form}
\citet[fn.~17]{Hausman1997} notes that the AIDS is a GGPF, so Gorman's theorem on exact two-stage
budgeting applies. The verification is short, and it shows which top level the theorem needs. Within group
$G$ the AIDS indirect utility is
\begin{equation*}
  u_G = \frac{\ln X_G - \ln a_G(p_G)}{B_G(p_G)}, \qquad B_G(p_G) = \beta_0\prod_{k\in G} p_k^{\beta_k}.
\end{equation*}
Here $\ln a_G = \ln\mathbf{P}_G$ is the translog index of the group. The function $B_G$ is homogeneous
of degree zero because the $\beta_k$ of the group sum to zero. Where $X_G > a_G(p_G)$, that is $u_G > 0$, the
increasing transformation $\ln u_G$ represents the same within-group preferences and equals
\begin{equation*}
  \ln u_G = \ln\ln\frac{X_G}{a_G(p_G)} - \ln B_G(p_G).
\end{equation*}
This is \eqref{eq:05-ggpf} with $F_G(z) = \ln\ln z$, $\mathcal{P}_G = a_G$ and $\mathcal{A}_G = -\ln B_G$.
Part (ii) of \cref{prop:05-first-stage} therefore applies when the top level is additive in the
transformed subutilities, $U = F\bigl(\sum_G \ln u_G\bigr)$. Under this assumption the exact group price
index of an AIDS group is its translog index $\mathbf{P}_G$, the index that
\citet{ChaudhuriGoldbergJia2006} use for the quinolone segment.

A separate increasing transformation of each subutility leaves within-group demands unchanged, but in general it
changes preferences across groups. Additivity in the untransformed $u_G$ is therefore not enough. Suppose
$U = \sum_G u_G$. The first-order conditions $\frac{\partial u_G}{\partial X_G} = \frac{1}{B_G X_G} =
\lambda$ give
\begin{equation*}
  X_G = I\,\frac{B_G(p_G)^{-1}}{\sum_H B_H(p_H)^{-1}} .
\end{equation*}
Group expenditure then varies with $B_G(p_G)$ even when $a_G(p_G)$ is held fixed. In the coordinates
$v_G = \ln u_G$ this aggregator is $\sum_G e^{v_G}$, which is not a function of $\sum_G v_G$.

\casehead{The first stage in practice}
Exact two-stage budgeting requires additivity at the top, which is itself a strong restriction on
substitution between groups. Applied work therefore treats the top level as an approximation. The
approximation of \citet[\S 5.2, eqs.~2.18--2.19, pp.~131--132]{DeatonMuellbauer1980b} starts from an
identity. Choose base prices $p^0$ with group subvectors $p_G^0$, and write each group expenditure function
as
\begin{equation}\label{eq:05-approx-stage}
  e_G(p_G, u_G) = e_G(p_G^0, u_G)\cdot\frac{e_G(p_G, u_G)}{e_G(p_G^0, u_G)}
                \equiv V_G\cdot P_G(p_G, p_G^0; u_G).
\end{equation}
The first factor $V_G$ is the money cost of reaching $u_G$ at base prices, a group quantity index
measured in money. The second factor is the exact cost-of-living index of the group at utility $u_G$, in
the sense of \cref{sec:05-cv-ev}. Under weak separability the first stage chooses subutilities to maximize
$F(u_1, \dots, u_{\mathcal{G}})$ subject to $\sum_G e_G(p_G, u_G) = I$. Since $V_G$ is increasing in
$u_G$, we can invert it and write $u_G = \chi_G(V_G)$, with base prices fixed. The first stage becomes
\begin{equation}\label{eq:05-approx-budget}
  \max_{V_1, \dots, V_{\mathcal{G}}} F\bigl(\chi_1(V_1), \dots, \chi_{\mathcal{G}}(V_{\mathcal{G}})\bigr)
  \quad\text{subject to}\quad \sum_G V_G\, P_G(p_G, p_G^0; u_G) = I .
\end{equation}
This looks like a standard consumer problem with quantities $V_G$ and prices $P_G$. It is not one, because
each $P_G$ depends on $u_G$, and allowing for this dependence takes us back to the original problem. The
dependence reflects within-group Engel curves that are not straight lines through the origin. If
homothetic subutilities remove it, \eqref{eq:05-approx-budget} is exact, which is route (i) of
\cref{prop:05-first-stage}. Otherwise the approximation replaces $P_G$ by an index that does not depend on
utility. When the variation of $P_G$ with $u_G$ is small, Laspeyres or Paasche indices of group prices
approximate it well. These are the indices that national accounts supply, so the construction
justifies standard practice, ``albeit in an approximate way'' (p.~132).

\citet[p.~13]{ChaudhuriGoldbergJia2006} adopt ``an approximate solution'' along these lines. It is an AIDS
for the shares of sub-segments in total expenditure, with group price indices in place of prices.
\citet{HausmanLeonardZona1994} also put group price indices in place of prices at their top and middle
levels (\cref{sec:05-hlz-model}). Hausman reports that he also estimated a
middle level that is not necessarily consistent with exact two-stage budgeting, with ``quite similar''
results \citep[fn.~17]{Hausman1997}. Group price indices may be the exact translog index of the lower
level, or a Stone or Laspeyres-type index of the group's prices (both defined in \cref{sec:05-stone}).
Hausman finds that this choice ``does
not typically have much influence on the final model estimates'' \citep[p.~218]{Hausman1997}.
Representative applications of multi-stage budgeting to differentiated products include
\citet{HausmanLeonardZona1994}, \citet{EllisonEtAl1997} and \citet{HausmanLeonard2002}.

\begin{remark}[Implicit separability]\label{rem:05-implicit-sep}
Weak separability is not the only structure that supports budgeting in stages. Preferences are
\emph{implicitly separable} if the expenditure function can be written as
\begin{equation*}
  e(p, u) = C\bigl[u, e_1(u, p_1), \dots, e_{\mathcal{G}}(u, p_{\mathcal{G}})\bigr],
\end{equation*}
where each $e_G$ is increasing in $u$ and linearly homogeneous in $p_G$ \citep{Gorman1970}. Total utility
$u$ enters every $e_G$, and there are no group subutilities. Shephard's lemma and the chain rule give
$q_i = \frac{\partial C}{\partial e_G}\frac{\partial e_G}{\partial p_i}$ for $i \in G$. Summing over the
group and using Euler's theorem for $e_G$,
\begin{equation*}
  X_G = \sum_{i \in G} p_iq_i = \frac{\partial C}{\partial e_G}\sum_{i \in G} p_i\frac{\partial e_G}{\partial p_i}
      = \frac{\partial C}{\partial e_G}\, e_G(u, p_G).
\end{equation*}
Dividing by $I = C$ and by $X_G$ gives the two stages:
\begin{equation*}
  \frac{X_G}{I} = \frac{\partial\ln C}{\partial\ln e_G}, \qquad
  w_{i|G} = \frac{\partial\ln e_G(u, p_G)}{\partial\ln p_i}.
\end{equation*}
The top stage allocates the budget with the ``macro'' cost function $C$ and the group price indices $e_G(u,
p_G)$. The bottom stage allocates within groups with the same indices. Both stages are exact, at the cost
of group price indices that depend on the level of total utility \citep[\S 5.2, eqs.~2.24--2.28,
pp.~133--134]{DeatonMuellbauer1980b}.
\end{remark}

\subsection{What separability implies for substitution}\label{sec:05-gu}

Separability buys a large reduction in parameters by restricting substitution between groups. The next
result makes the restriction precise.

\begin{proposition}[\citealp{GoldmanUzawa1964}]\label{prop:05-goldman-uzawa}
Suppose preferences are weakly separable and demands are differentiable, with nonzero income effects
$\frac{\partial X_G}{\partial I}$. For goods $i \in G$ and $j \in G'$ in different groups, the Slutsky
terms satisfy
\begin{equation}\label{eq:05-gu}
  S_{ij} = \mu_{GG'}(p, I)\, \frac{\partial q_i}{\partial I}\, \frac{\partial q_j}{\partial I}
\end{equation}
for a function $\mu_{GG'}$ that depends on the pair of groups but not on the goods within them.
\end{proposition}

\begin{proof}
We follow the argument of \citet{Gorman1971} as presented by \citet[\S 5.2, eqs.~2.4--2.9,
pp.~128--129]{DeatonMuellbauer1980b}. By \cref{prop:05-second-stage}, $q_i(p, I) = q_{i|G}(p_G, X_G(p, I))$. The Hicksian demand is $h_i(p, u) =
q_i(p, e(p, u)) = q_{i|G}(p_G, X_G^h(p, u))$, where $X_G^h(p,u) = X_G(p, e(p,u))$ is compensated group
expenditure. For $j \notin G$, the price $p_j$ enters $h_i$ only through $X_G^h$, so
\begin{equation*}
  S_{ij} = \frac{\partial h_i}{\partial p_j}
         = \frac{\partial q_{i|G}}{\partial X_G}\,\frac{\partial X_G^h}{\partial p_j}.
\end{equation*}
The income derivative has the same structure, $\frac{\partial q_i}{\partial I} = \frac{\partial
q_{i|G}}{\partial X_G}\frac{\partial X_G}{\partial I}$. Eliminating $\frac{\partial q_{i|G}}{\partial
X_G}$,
\begin{equation*}
  S_{ij} = \frac{\partial q_i}{\partial I}\left(\frac{\partial X_G}{\partial I}\right)^{-1}
           \frac{\partial X_G^h}{\partial p_j}.
\end{equation*}
The same steps with the roles of $i$ and $j$ exchanged give $S_{ji} = \frac{\partial q_j}{\partial I}
(\frac{\partial X_{G'}}{\partial I})^{-1}\frac{\partial X_{G'}^h}{\partial p_i}$. Slutsky symmetry,
$S_{ij} = S_{ji}$ (\cref{prop:05-slutsky}), and division by $\frac{\partial q_i}{\partial
I}\frac{\partial q_j}{\partial I}$ give
\begin{equation*}
  \left(\frac{\partial q_j}{\partial I}\right)^{-1}\frac{\partial X_G^h}{\partial p_j}
  \left(\frac{\partial X_G}{\partial I}\right)^{-1}
  = \left(\frac{\partial q_i}{\partial I}\right)^{-1}\frac{\partial X_{G'}^h}{\partial p_i}
  \left(\frac{\partial X_{G'}}{\partial I}\right)^{-1}.
\end{equation*}
The left side does not depend on which good $i \in G$ we picked, and the right side does not depend on
which good $j \in G'$ we picked. Since the equality holds for all such pairs, both sides equal a common
value $\mu_{GG'}(p, I)$. The left side then gives $\frac{\partial X_G^h}{\partial
p_j}\bigl(\frac{\partial X_G}{\partial I}\bigr)^{-1} = \mu_{GG'}\frac{\partial q_j}{\partial I}$.
Substituting into the expression for $S_{ij}$,
\begin{equation*}
  S_{ij} = \mu_{GG'}\,\frac{\partial q_i}{\partial I}\,\frac{\partial q_j}{\partial I},
\end{equation*}
which is \eqref{eq:05-gu}. If a good has a zero income effect the formula holds trivially for the pairs
that involve it.
\end{proof}

\citet{Gorman1971} shows that the restriction \eqref{eq:05-gu} is also sufficient for weak separability,
so it summarizes the empirical content of separability for substitution
\citep[p.~129]{DeatonMuellbauer1980b}. The factor $\mu_{GG'}$ has an
interpretation in terms of the groups themselves. Define
\begin{equation}\label{eq:05-lambda-GG}
  \lambda_{GG'} = \mu_{GG'}\,\frac{\partial X_G}{\partial I}\,\frac{\partial X_{G'}}{\partial I}.
\end{equation}
Since $\frac{\partial q_i}{\partial I} = \frac{\partial q_{i|G}}{\partial X_G}\frac{\partial X_G}{\partial
I}$, \eqref{eq:05-gu} becomes $S_{ij} = \lambda_{GG'}\frac{\partial q_{i|G}}{\partial
X_G}\frac{\partial q_{j|G'}}{\partial X_{G'}}$. The first display of the proof gives $S_{ij} =
\frac{\partial q_{i|G}}{\partial X_G}\frac{\partial X_G^h}{\partial p_j}$. Equating the two expressions,
\begin{equation}\label{eq:05-XG-pj}
  \frac{\partial X_G^h}{\partial p_j} = \lambda_{GG'}\,\frac{\partial q_{j|G'}}{\partial X_{G'}},
  \qquad j \in G',\ G' \neq G .
\end{equation}
Multiply by $p_j$ and sum over $j \in G'$. Conditional demands add up to group expenditure, so $\sum_{j \in
G'} p_j\frac{\partial q_{j|G'}}{\partial X_{G'}} = 1$, and
\begin{equation}\label{eq:05-lambda-interp}
  \lambda_{GG'} = \sum_{j \in G'} p_j\,\frac{\partial X_G^h}{\partial p_j}
\end{equation}
\citep[\S 5.2, eq.~2.10, p.~129]{DeatonMuellbauer1980b}. The right side is the compensated response of
spending on group $G$ to a proportional rise in all prices of group $G'$. So $\lambda_{GG'}$ is the
substitution term between two group aggregates, each defined with fixed relative prices inside the group.
Writing $X_G^h = \sum_{i \in G} p_ih_i$ in \eqref{eq:05-lambda-interp} gives $\lambda_{GG'} = \sum_{i \in G}
\sum_{j \in G'} p_ip_jS_{ij}$, which is symmetric in $(G, G')$ by Slutsky symmetry. We define $\lambda_{GG}$
by the same double sum with $G' = G$. Let $\Pi$ be the $J \times \mathcal{G}$ matrix with elements $p_j$
in the rows of the goods of group $G$ and column $G$, and zeros elsewhere. Then $\Lambda = [\lambda_{GG'}]
= \Pi^{\T} S\,\Pi$. This matrix is symmetric and negative semidefinite. Its rows sum to zero because
$\Pi\iota = p$ for the vector of ones $\iota$, and $Sp = 0$. It is the Slutsky matrix of the group
aggregates. Weak separability thus gives
two tiers of substitution matrices. Within each group the conditional Slutsky matrix is unrestricted, and
between groups substitution runs through $\Lambda$ alone.

The signs of cross-group substitution are also restricted. If no good is inferior, $\frac{\partial
q_i}{\partial I} > 0$ for all $i$, and \eqref{eq:05-gu} gives every $S_{ij}$ with $i \in G$ and $j \in G'$
the sign of $\mu_{GG'}$. When two groups are substitutes, every good of one is a Hicksian substitute for
every good of the other. When the groups are complements, every cross pair is complementary
\citep[p.~129]{DeatonMuellbauer1980b}.

In elasticity form, with $\frac{\partial q_i}{\partial I} = \eta_i\frac{q_i}{I}$, \eqref{eq:05-gu} becomes
\begin{equation}\label{eq:05-gu-elast}
  \varepsilon^h_{ij} = \frac{\mu_{GG'}}{I}\,\eta_i\,\eta_j\, w_j, \qquad i \in G,\ j \in G',\ G \neq G'.
\end{equation}
All compensated cross elasticities between two groups are governed by one number per pair of groups,
scaled by the goods' expenditure elasticities and the budget share of the good whose price changes. Under
additive separability the factor $\mu$ is the same for all pairs of groups. Substitution between groups
runs entirely through group expenditures.

The restriction is strong when goods in different groups are close substitutes. If a cheap good in one
group competes directly with a particular good in another, separability spreads that competition evenly,
in proportion to income effects, over all goods of the second group. The sign result adds a second
constraint. Two groups that are substitutes on the whole cannot contain a complementary cross pair, such
as a camera in one group and its film in the other. The grouping must therefore put close substitutes
together and keep complements in the same group. In ready-to-eat cereals the segmentation comes from marketing research and
has little ambiguity, but \citet[fn.~5]{Bresnahan1997} warns that this is an exception. In many
industries ``several principles of differentiation, each with a distinct set of close-substitute
products, compete for the analyst's attention.'' In automobiles, size segments, brand names and even
country reputations all shape substitution, and many consumer product markets mix named segments with a
private-label versus brand-name split.

Separability can be tested at two levels. Under weak separability, the prices of goods outside a group
affect the conditional demands within the group only through group expenditure $X_G$. Adding individual
outside prices to a within-group equation that conditions on observed $X_G$ and testing their exclusion
is therefore a test of the grouping. A group-level equation explains the demand for a group by one price
index per group. Outside prices then enter only through the indices if group expenditures depend on
prices only through them, as under the conditions of \cref{prop:05-first-stage}. Weak separability alone
does not deliver this (\cref{sec:05-first-stage}). An exclusion test in a group-level equation is
therefore a joint test of the grouping and of this index aggregation. \citet{HausmanLeonardZona1994}
implement such a test by adding the prices of individual light beers to the segment equation for premium
beer (\cref{sec:05-hlz}).

Under weak separability alone, the first stage takes a differential form with two price indices per
group. The result goes back to \citet{BartenTurnovsky1966}, and we
state it as in \citet[\S 5.2, eqs.~2.20--2.23, pp.~132--133]{DeatonMuellbauer1980b}. Write $w_i =
\frac{p_iq_i}{I}$ for the share of good $i$ in total expenditure, $w_G = \frac{X_G}{I}$ for the share of
group $G$, and $b_G = \sum_{i \in G} p_i\frac{\partial q_i}{\partial I}$ for the marginal propensity to
spend on the group. Define two differential price indices for each group,
\begin{equation}\label{eq:05-two-indices}
  \dd\ln P^{w}_G = \sum_{k \in G} w_{k|G}\,\dd\ln p_k, \qquad
  \dd\ln P^{m}_G = \sum_{k \in G} p_k\frac{\partial q_{k|G}}{\partial X_G}\,\dd\ln p_k .
\end{equation}
The first weights price changes by within-group budget shares. The second weights them by within-group
marginal budget shares. Both sets of weights sum to one.

\begin{proposition}[Two price indices per group]\label{prop:05-two-indices}
Under weak separability, for every group $G$,
\begin{equation}\label{eq:05-bt}
  \sum_{i \in G} w_i\,\dd\ln q_i
  = b_G\Bigl(\dd\ln I - \sum_{G'} w_{G'}\,\dd\ln P^{w}_{G'}\Bigr)
  + \sum_{G'} \frac{\lambda_{GG'}}{I}\,\dd\ln P^{m}_{G'} ,
\end{equation}
where the sums run over all groups, including $G$, and $\lambda_{GG'}$ is defined in
\eqref{eq:05-lambda-GG} and below \eqref{eq:05-lambda-interp}.
\end{proposition}

\begin{proof}
Totally differentiate $q_i(p, I)$ and substitute the Slutsky equation $\frac{\partial q_i}{\partial p_j} =
S_{ij} - q_j\frac{\partial q_i}{\partial I}$ (\cref{prop:05-slutsky}):
\begin{equation*}
  \dd q_i = \frac{\partial q_i}{\partial I}\Bigl(\dd I - \sum_j q_j\,\dd p_j\Bigr) + \sum_j S_{ij}\,\dd p_j .
\end{equation*}
Multiply by $\frac{p_i}{I}$, use $\dd p_j = p_j\,\dd\ln p_j$, and sum over $i \in G$:
\begin{equation*}
  \sum_{i \in G} w_i\,\dd\ln q_i = b_G\Bigl(\dd\ln I - \sum_j w_j\,\dd\ln p_j\Bigr)
  + \frac{1}{I}\sum_j\Bigl(\sum_{i \in G} p_iS_{ij}\Bigr)p_j\,\dd\ln p_j .
\end{equation*}
Since $w_j = w_{G'}w_{j|G'}$ for $j \in G'$, the first bracket equals $\dd\ln I - \sum_{G'}
w_{G'}\,\dd\ln P^{w}_{G'}$. It remains to show that, for every $j$ in any group $G'$,
\begin{equation}\label{eq:05-bt-step}
  \sum_{i \in G} p_iS_{ij} = \lambda_{GG'}\,\frac{\partial q_{j|G'}}{\partial X_{G'}} .
\end{equation}
For $G' \neq G$, the sum equals $\frac{\partial X_G^h}{\partial p_j}$, and \eqref{eq:05-XG-pj} gives
\eqref{eq:05-bt-step}. For $G' = G$, use Slutsky symmetry and then homogeneity of $h_j$ of degree zero in
prices, $\sum_k p_kS_{jk} = 0$:
\begin{equation*}
  \sum_{i \in G} p_iS_{ij} = \sum_{i \in G} p_iS_{ji} = -\sum_{G'' \neq G}\ \sum_{k \in G''} p_kS_{kj}
  = -\frac{\partial q_{j|G}}{\partial X_G}\sum_{G'' \neq G}\lambda_{G''G} .
\end{equation*}
The last step applies \eqref{eq:05-XG-pj} to group $G''$ and the price $p_j$ of a good in $G$. The rows
of $\Lambda$ sum to zero and $\Lambda$ is symmetric, so $-\sum_{G'' \neq G}\lambda_{G''G} = \lambda_{GG}$,
which proves \eqref{eq:05-bt-step} for $G' = G$. Substituting \eqref{eq:05-bt-step} into the last sum,
\begin{equation*}
  \frac{1}{I}\sum_{G'}\lambda_{GG'}\sum_{j \in G'} p_j\frac{\partial q_{j|G'}}{\partial X_{G'}}\,\dd\ln p_j
  = \sum_{G'}\frac{\lambda_{GG'}}{I}\,\dd\ln P^{m}_{G'} ,
\end{equation*}
which completes the proof.
\end{proof}

The left side of \eqref{eq:05-bt} is the change in a share-weighted quantity index of group $G$. The
right side has the form of a Rotterdam equation for groups (\cref{sec:05-rotterdam}). Real income enters
through the share-weighted indices $P^w$, and substitution through the marginally weighted indices $P^m$.
Knowledge of two price indices per group is therefore enough to maintain an optimal allocation across
groups, once it is reached. No restriction beyond weak separability is needed. In practice, if prices
inside each group move roughly together, the two indices move together and one index may be an adequate
proxy for both (p.~133). The result is local. The coefficients $b_G$ and $\lambda_{GG'}$ vary with prices and
expenditure, and a Rotterdam-type model treats them as constants.

\subsection{Conditional and unconditional elasticities}\label{sec:05-uncond}

Estimates of a multi-level system yield \emph{conditional} elasticities, which hold group expenditure
fixed. Pricing, merger and welfare calculations need \emph{unconditional} elasticities, which let group
expenditure respond. As elsewhere in this chapter, elasticities are signed: $\varepsilon_{ij} =
\frac{\partial\ln q_i}{\partial\ln p_j}$, so own-price elasticities are negative.

Consider a two-level system. Total expenditure on the upper-level category is $X$, for example spending
on all systemic antibiotics, and plays the role of income. The upper level allocates $X$ to groups with
shares $w_G = \frac{X_G}{X}$ that depend on group price indices $P_{G'}$ and on $X$. The lower level is a
conditional demand system $q_{i|G}(p_G, X_G)$ within each group, and $P_G = P_G(p_G)$. For $i \in G$ and
a price $p_j$ of the same group, the chain rule with $X$ held fixed gives
\begin{equation}\label{eq:05-uncond-elast}
  \varepsilon_{ij} = \varepsilon_{ij}\big|_{X_G}
  + \frac{\partial\ln q_i}{\partial\ln X_G}\,\frac{\partial\ln X_G}{\partial\ln P_G}\,
    \frac{\partial\ln P_G}{\partial\ln p_j},
\end{equation}
which is equation (3.10) of \citet{ChaudhuriGoldbergJia2006}. The first term is the conditional
elasticity. The second term traces the effect of $p_j$ on the group price index, of the index on group
expenditure, and of group expenditure on $q_i$.

With an AIDS at both levels, each factor has a closed form. The lower level gives the conditional
Marshallian elasticity of \eqref{eq:05-elast-marshall}, written with within-group shares,
\begin{equation*}
  \varepsilon_{ij}\big|_{X_G} = -\indic(i = j)
  + \frac{\gamma_{ij} - \beta_i\bigl(w_{j|G} - \beta_j\ln\frac{X_G}{P_G}\bigr)}{w_{i|G}},
\end{equation*}
and the conditional expenditure elasticity $\frac{\partial\ln q_i}{\partial\ln X_G} = 1 +
\frac{\beta_i}{w_{i|G}}$. With the translog group index, Shephard's lemma at the lower level gives
\begin{equation*}
  \frac{\partial\ln P_G}{\partial\ln p_j} = \alpha_j + \sum_{k\in G}\gamma_{jk}\ln p_k
  = w_{j|G} - \beta_j\ln\frac{X_G}{P_G},
\end{equation*}
which reduces to $w_{j|G}$ for a Stone group index whose weights are held fixed. At the upper level,
$X_G = w_G X$ with $X$ fixed, so
\begin{equation}\label{eq:05-dlnXG}
  \frac{\partial\ln X_G}{\partial\ln P_G} = \frac{\partial\ln w_G}{\partial\ln P_G}
  = \frac{1}{w_G}\frac{\partial w_G}{\partial\ln P_G} = 1 + \varepsilon_{GG},
\end{equation}
where $\varepsilon_{GG}$ is the own-price elasticity of the group quantity index $\frac{X_G}{P_G}$. With
an upper-level AIDS deflated by a Stone index, $\ln P = \sum_{G'} w_{G'}\ln P_{G'}$, and the Stone weights
held fixed, $\frac{\partial w_G}{\partial\ln P_G} = \gamma_{GG} - \beta_G w_G$. Collecting terms,
\begin{equation}\label{eq:05-uncond-aids}
  \varepsilon_{ij} = \varepsilon_{ij}\big|_{X_G}
  + \Bigl(1 + \frac{\beta_i}{w_{i|G}}\Bigr)\Bigl(\frac{\gamma_{GG}}{w_G} - \beta_G\Bigr)
    \Bigl(w_{j|G} - \beta_j\ln\frac{X_G}{P_G}\Bigr).
\end{equation}
The middle factor shows when the correction matters. If the group as a whole has unit-elastic demand
($\varepsilon_{GG} = -1$), a higher group price leaves group expenditure unchanged and unconditional and
conditional elasticities coincide. If group demand is inelastic, a higher group price draws spending into
the group. If it is elastic, spending leaves the group. The sign of the correction is the product of the
signs of the three factors in \eqref{eq:05-uncond-elast}. Suppose good $i$ is normal within its group, $1
+ \frac{\beta_i}{w_{i|G}} > 0$, and $p_j$ raises the group index, $w_{j|G} - \beta_j\ln\frac{X_G}{P_G} >
0$. The correction is then positive when group demand is inelastic and negative when it is elastic. Both
signs reverse for a good that is inferior within its group. In the AIDS this happens when $\beta_i < 0$
and $w_{i|G} < -\beta_i$, and a negative semidefinite Slutsky matrix does not rule it out.

For goods in different groups, $i \in G$ and $j \in G'$ with $G' \neq G$, the conditional term vanishes
and only the expenditure channel remains:
\begin{equation}\label{eq:05-uncond-cross}
  \varepsilon_{ij} = \frac{\partial\ln q_i}{\partial\ln X_G}\,\frac{\partial\ln X_G}{\partial\ln P_{G'}}\,
  \frac{\partial\ln P_{G'}}{\partial\ln p_j}
  = \Bigl(1 + \frac{\beta_i}{w_{i|G}}\Bigr)\,\varepsilon_{GG'}\,
    \Bigl(w_{j|G'} - \beta_j\ln\frac{X_{G'}}{P_{G'}}\Bigr),
\end{equation}
where $\varepsilon_{GG'} = \frac{1}{w_G}\frac{\partial w_G}{\partial\ln P_{G'}} = \frac{\gamma_{GG'} -
\beta_G w_{G'}}{w_G}$ is the upper-level cross elasticity. This is the elasticity counterpart of
\cref{prop:05-goldman-uzawa}: cross-group effects are the product of a group-level elasticity and
within-group factors. Finally, the unconditional elasticity of $q_i$ with respect to the category budget
is the product of the expenditure elasticities at the two levels,
\begin{equation*}
  \frac{\partial\ln q_i}{\partial\ln X} = \frac{\partial\ln q_i}{\partial\ln X_G}\,
  \frac{\partial\ln X_G}{\partial\ln X} = \Bigl(1 + \frac{\beta_i}{w_{i|G}}\Bigr)
  \Bigl(1 + \frac{\beta_G}{w_G}\Bigr).
\end{equation*}

\casehead{Three levels}
\citet{HausmanLeonardZona1994} and \citet{Hausman1997} use three levels. The bottom level explains brand
shares within a segment by an AIDS on segment expenditure. The middle level explains segment quantities
by a log-log system in category expenditure and segment price indices. The top level explains category
demand by a log-log equation in income and the category price index. Let
$\Pi$ denote the category price index and $y$ income. Group expenditure $X_G(P_1, \dots,
P_{\mathcal{G}}, X)$ now responds to $P_G$ also through the category budget $X(\Pi, y)$. The total effect
of the segment index on segment expenditure is
\begin{equation*}
  \frac{\dd\ln X_G}{\dd\ln P_G} = \frac{\partial\ln X_G}{\partial\ln P_G}\Big|_{X}
  + \frac{\partial\ln X_G}{\partial\ln X}\,\frac{\partial\ln X}{\partial\ln\Pi}\,
    \frac{\partial\ln\Pi}{\partial\ln P_G},
\end{equation*}
and this total replaces $\frac{\partial\ln X_G}{\partial\ln P_G}$ in \eqref{eq:05-uncond-elast}. The
estimating equations of the three levels are given in \cref{sec:05-hlz}.

\Cref{fig:05-budget-tree} draws the two budgeting trees used in this chapter.

\begin{figure}[htbp]
\centering
\begin{tikzpicture}[
  lvl/.style={draw, rounded corners=2pt, font=\scriptsize, align=center, inner sep=3pt,
    minimum height=0.62cm},
  edge/.style={-{Stealth[length=1.6mm]}, thin},
  note/.style={font=\scriptsize\itshape, align=left, text width=3.3cm, anchor=west}]
\node[lvl] (inc) at (4.6,0) {Income $y$};
\node[lvl] (cat) at (2.6,-1.25) {Cereal, $X$ (index $\Pi$)};
\node[lvl] (oth) at (6.8,-1.25) {Other goods};
\node[lvl] (ad) at (0.0,-2.5) {Adult, $X_1$};
\node[lvl] (ch) at (2.6,-2.5) {Child, $X_2$};
\node[lvl] (fa) at (5.2,-2.5) {Family, $X_3$};
\node[lvl] (b1) at (3.4,-3.75) {Cheerios};
\node[lvl] (b2) at (5.2,-3.75) {Corn Flakes};
\node[lvl, dashed] (b3) at (7.1,-3.75) {other brands};
\draw[edge] (inc) -- (cat); \draw[edge] (inc) -- (oth);
\draw[edge] (cat) -- (ad); \draw[edge] (cat) -- (ch); \draw[edge] (cat) -- (fa);
\draw[edge] (fa) -- (b1); \draw[edge] (fa) -- (b2); \draw[edge] (fa) -- (b3);
\node[note] at (9.0,-1.25) {top: category demand on income and $\Pi$};
\node[note] at (9.0,-2.5) {middle: segment demand on $X$ and segment indices};
\node[note] at (9.0,-3.75) {bottom: brand shares on segment expenditure};
\node[font=\small] at (-1.3,0) {(a)};
\end{tikzpicture}

\vspace{0.6cm}
\begin{tikzpicture}[
  lvl/.style={draw, rounded corners=2pt, font=\scriptsize, align=center, inner sep=2pt,
    minimum height=0.62cm},
  seg/.style={draw, rounded corners=2pt, font=\scriptsize, align=center, inner sep=2pt,
    text width=1.42cm, minimum height=0.8cm},
  grp/.style={draw, rounded corners=2pt, font=\scriptsize, align=center, inner sep=2pt,
    text width=1.66cm, minimum height=0.8cm},
  edge/.style={-{Stealth[length=1.6mm]}, thin}]
\node[lvl] (top) at (6.3,0) {Systemic antibacterials, total expenditure $X$};
\foreach \name/\x [count=\n] in {Tetra-\\cyclines/0.0, Chloram-\\phenicol/1.8, Ampicillin/3.6,
    Cephalo-\\sporins/5.4, Trimetho-\\prim/7.2, Macrolides/9.0, Other\\penicillins/10.8,
    Quinolones/12.6} {
  \node[seg] (s\n) at (\x,-1.45) {\name};
  \draw[edge] (top.south) -- (s\n.north);
}
\foreach \name/\x [count=\n] in {Foreign\\ciprofloxacin/0.6, Foreign\\norfloxacin/2.6,
    Foreign\\ofloxacin/4.6, Domestic\\ciprofloxacin/6.6, Domestic\\norfloxacin/8.6,
    Domestic\\ofloxacin/10.6, Domestic\\sparfloxacin/12.6} {
  \node[grp] (g\n) at (\x,-3.3) {\name};
  \draw[edge] (s8.south) -- (g\n.north);
}
\node[lvl, dashed] (sku1) at (6.6,-4.8) {SKUs (firm $\times$ presentation)};
\node[lvl, dashed] (sku2) at (10.6,-4.8) {SKUs (firm $\times$ presentation)};
\draw[edge, dashed] (g4.south) -- (sku1.north);
\draw[edge, dashed] (g6.south) -- (sku2.north);
\node[font=\small] at (-1.0,0) {(b)};
\end{tikzpicture}
\caption{Budgeting trees. Panel (a): the three-level structure of \citet{Hausman1997} for ready-to-eat
cereal, with the three segments he uses and two of the family-segment brands (his Table 5.1). Panel (b):
the two-level structure of \citet{ChaudhuriGoldbergJia2006}. The upper level allocates systemic
antibacterial expenditure across sub-segments, eight of which appear with equations in their
upper-level system (their Table A3). The lower level allocates quinolone expenditure across seven
molecule-by-status groups. Foreign sparfloxacin, with a minuscule share, is folded into the domestic
group. Within-group choice among stock-keeping units (SKUs) is not modeled (dashed).}
\label{fig:05-budget-tree}
\end{figure}

\section{Estimation and identification}\label{sec:05-estimation}

This section turns the AIDS into an econometric model. We add an error term and ask what it contains,
deal with the singular covariance matrix that adding up creates, and compare the exact nonlinear model
with its linear approximation. We then show how the theoretical restrictions are tested and work through
the original application of \citet{DeatonMuellbauer1980a} in detail. The section closes with the
identification problems that dominate industrial organization applications: endogenous prices,
endogenous group expenditure, and group prices measured with error.

\subsection{Stochastic specification}\label{sec:05-stochastic}

Observations are indexed by $t = 1, \dots, T$. An observation may be a year of national data, a month in a
region, or a week in a city. The estimating equations are
\begin{equation}\label{eq:05-stoch}
  w_{it} = \alpha_i + \sum_{j=1}^{J}\gamma_{ij}\ln p_{jt} + \beta_i\ln\frac{I_t}{\mathbf{P}_t} + \xi_{it},
  \qquad i = 1, \dots, J,
\end{equation}
with the translog index $\ln\mathbf{P}_t = \alpha_0 + \sum_k\alpha_k\ln p_{kt} + \frac{1}{2}\sum_k\sum_j
\gamma_{kj}\ln p_{kt}\ln p_{jt}$ of \eqref{eq:05-translog}. We write the error as $\xi_{it}$, in place of
the $\varepsilon_{it}$ used by the original papers, because $\varepsilon$ denotes elasticities in this
chapter. Its content matters for every later step, so we list the main components.
\begin{enumerate}[label=(\alph*)]
  \item \emph{Unobserved demand shifters.} Advertising campaigns, epidemics that call for a particular
  drug, or changes in demographics shift the shares of particular goods. Firms often observe these
  shocks when they set prices.
  \item \emph{Measurement error.} Shares and prices are measured with error. In multi-level systems the
  prices of product groups are averages of the prices of many stock-keeping units, which approximate an
  unobserved exact group price index (\cref{sec:05-endogeneity}).
  \item \emph{Aggregation error.} The omitted term $-\beta_i(\ln\kappa_t - \ln\kappa^*)$ of
  \eqref{eq:05-aids-agg-ts} enters the error when the distribution of budgets changes.
  \item \emph{Specification error.} Habits, stocks of durables, price expectations and other dynamics
  that the static model omits.
\end{enumerate}
Component (a) makes prices endogenous when firms set them. Component (b) creates errors in variables.
Components (c) and (d) are correlated with the regressors whenever inequality or the omitted dynamics move
with prices and real expenditure. We write $\Sigma = \E[\xi_t\xi_t^{\T}]$ for the $J\times J$ covariance
matrix of $\xi_t = (\xi_{1t}, \dots, \xi_{Jt})^{\T}$.

\subsection{Adding up and the singular covariance matrix}\label{sec:05-singular}

Budget shares sum to one in the data. Under the adding-up restrictions of \cref{prop:05-restrictions},
the deterministic part of \eqref{eq:05-stoch} also sums to one, because
\begin{equation*}
  \sum_i\Bigl(\alpha_i + \sum_j\gamma_{ij}\ln p_{jt} + \beta_i\ln\frac{I_t}{\mathbf{P}_t}\Bigr)
  = \sum_i\alpha_i + \sum_j\Bigl(\sum_i\gamma_{ij}\Bigr)\ln p_{jt} + \Bigl(\sum_i\beta_i\Bigr)
    \ln\frac{I_t}{\mathbf{P}_t} = 1 .
\end{equation*}
Hence $\sum_i\xi_{it} = 0$ for every $t$. Writing $\iota$ for the $J$-vector of ones, $\iota^{\T}\xi_t = 0$
implies $\iota^{\T}\Sigma = 0$, so $\Sigma$ is singular with rank at most $J - 1$. The Gaussian likelihood
of the full system is not defined, because it requires $\Sigma^{-1}$.

The standard solution drops one equation, say the $J$-th, and estimates the remaining $J - 1$. The
parameters of the dropped equation follow from adding up:
\begin{equation*}
  \alpha_J = 1 - \sum_{i<J}\alpha_i, \qquad \beta_J = -\sum_{i<J}\beta_i, \qquad
  \gamma_{Jj} = -\sum_{i<J}\gamma_{ij} \quad (j = 1, \dots, J).
\end{equation*}
Which equation we drop does not matter for maximum likelihood.

\begin{proposition}[\citealp{Barten1969}]\label{prop:05-barten}
Suppose the errors of any $J-1$ share equations are jointly normal with an unrestricted covariance matrix.
Then the maximum likelihood estimates of the parameters are the same whichever equation is dropped.
\end{proposition}

\begin{proof}
Let $\tilde\xi_t(\theta)$ be the residuals of the $J-1$ retained equations at parameter value $\theta$.
Concentrating the unrestricted covariance matrix out of the Gaussian log likelihood gives, up to a
constant,
\begin{equation*}
  \ln L^c(\theta) = -\frac{T}{2}\ln\det\hat\Omega(\theta), \qquad
  \hat\Omega(\theta) = \frac{1}{T}\sum_t\tilde\xi_t(\theta)\tilde\xi_t(\theta)^{\T}.
\end{equation*}
Dropping equation $k$ instead of equation $J$ replaces $\xi_{kt}$ by $\xi_{Jt} = -\sum_{i<J}\xi_{it}$. The
new residual vector is $A\tilde\xi_t$, where $A$ is the identity matrix with row $k$ replaced by a row of
$-1$'s, followed by a reordering of rows. Expanding the determinant along row $k$ shows that $\det A =
\pm 1$. The new concentrated likelihood is therefore
\begin{equation*}
  -\frac{T}{2}\ln\det\bigl(A\hat\Omega(\theta)A^{\T}\bigr)
  = -\frac{T}{2}\ln\bigl[(\det A)^2\det\hat\Omega(\theta)\bigr] = \ln L^c(\theta),
\end{equation*}
the same function of $\theta$, with the same maximizer.
\end{proof}

\casehead{Seemingly unrelated regressions}
The standard system estimator is \emph{seemingly unrelated regressions} (SUR) \citep{Zellner1962}. Write
retained equation $i$ as $\mathbf{w}_i = \mathbf{X}_ib_i + \boldsymbol{\xi}_i$, where $\mathbf{w}_i$ stacks
the $T$ observed shares, $\mathbf{X}_i$ is the $T \times K_i$ matrix of its regressors and $b_i$ is its
coefficient vector. Stacking the $J - 1$ equations gives $\mathbf{w} = \mathbf{X}b + \boldsymbol{\xi}$, with
$\mathbf{X} = \diag(\mathbf{X}_1, \dots, \mathbf{X}_{J-1})$ block diagonal. The errors of one observation
have covariance matrix $\E[\tilde\xi_t\tilde\xi_t^{\T}] = \Omega$ and are independent over $t$, so
$\Var(\boldsymbol{\xi}) = \Omega\otimes I_T$. Feasible generalized least squares gives
\begin{equation}\label{eq:05-sur}
  \hat b = \bigl[\mathbf{X}^{\T}(\hat\Omega^{-1}\otimes I_T)\mathbf{X}\bigr]^{-1}
           \mathbf{X}^{\T}(\hat\Omega^{-1}\otimes I_T)\,\mathbf{w},
\end{equation}
where $\hat\Omega$ has elements $\hat\Omega_{ik} = \frac{1}{T}\hat{\boldsymbol{\xi}}_i^{\T}\hat{\boldsymbol{\xi}}_k$
computed from equation-by-equation OLS residuals. Homogeneity and symmetry are linear restrictions
$Rb = r_0$. Minimizing the generalized least squares criterion subject to them gives
\begin{equation}\label{eq:05-sur-restricted}
  \hat b_R = \hat b - VR^{\T}\bigl(RVR^{\T}\bigr)^{-1}\bigl(R\hat b - r_0\bigr), \qquad
  V = \bigl[\mathbf{X}^{\T}(\hat\Omega^{-1}\otimes I_T)\mathbf{X}\bigr]^{-1}.
\end{equation}
\Cref{sec:A-systems} derives this formula from the Lagrangian. Recomputing $\hat\Omega$ from the new residuals
and iterating never lowers the Gaussian likelihood. A limit of the iteration is a stationary point of that
likelihood, typically a local maximum. It is the maximum likelihood estimator only if it is the global maximum
(\cref{sec:A-systems}). Restarting from several points guards against a local maximum (\cref{sec:A-numerical}).
The proof of \cref{prop:05-barten} shows that the concentrated likelihood is the same function of the parameters
whichever equation is dropped. Every stationary point that iterated SUR can reach therefore shares that
invariance. \emph{Three-stage least squares} applies \eqref{eq:05-sur} after
replacing $\mathbf{X}$ by its projection on the instruments, with $\hat\Omega$ estimated from two-stage least
squares residuals (\cref{sec:A-3sls}).

Two further facts simplify estimation when there are no cross-equation restrictions. First, when every
equation has the same regressors, $\mathbf{X}_i = \mathbf{X}_0$ for all $i$, SUR reduces to ordinary least
squares (OLS) equation by equation (\cref{prop:A-sur-ols}). OLS is then the maximum likelihood estimator
under normality \citep[p.~316]{DeatonMuellbauer1980a}. Second, OLS automatically satisfies adding up. To see
this, stack the $T$ observations of the regressors $(1, \ln p_{1t}, \dots, \ln p_{Jt},
\ln\frac{I_t}{\mathbf{P}_t})$ in the matrix $\mathbf{X}_0$, whose first column is the vector of ones $\iota_T$.
The OLS coefficient vectors $b_i = (\mathbf{X}_0^{\T}\mathbf{X}_0)^{-1}\mathbf{X}_0^{\T}\mathbf{w}_i$ sum to
\begin{equation*}
  \sum_i b_i = (\mathbf{X}_0^{\T}\mathbf{X}_0)^{-1}\mathbf{X}_0^{\T}\sum_i \mathbf{w}_i
  = (\mathbf{X}_0^{\T}\mathbf{X}_0)^{-1}\mathbf{X}_0^{\T}\iota_T = (1, 0, \dots, 0)^{\T},
\end{equation*}
because regressing a column of ones on $\mathbf{X}_0$ returns a coefficient of one on the constant and zero
elsewhere. Hence $\sum_i\hat\alpha_i = 1$, $\sum_i\hat\gamma_{ij} = 0$ and $\sum_i\hat\beta_i = 0$.

\casehead{Short samples}
Symmetry is a cross-equation restriction, so the covariance matrix now affects the estimates. With few
observations an unrestricted covariance matrix cannot be estimated. \citet[p.~318]{DeatonMuellbauer1980a}
had 21 annual observations on eight shares. They note that the likelihood ``can be made arbitrarily
large by making any one equation fit perfectly.'' The reason is visible in $\ln L^c(\theta)$. If the
parameters can be chosen so that one equation has zero residuals, $\det\hat\Omega(\theta) = 0$ and the
concentrated likelihood is unbounded. Under symmetry their system has 42 free coefficients, and every
equation depends on all of them through $\ln\mathbf{P}_t$. With only 21 observations per equation, the
coefficients can typically be chosen to fit one equation exactly. Deaton and Muellbauer therefore impose
a structure on the covariance matrix,
\begin{equation}\label{eq:05-dm-cov}
  \Sigma = \sigma^2\Bigl(I_J - \frac{1}{J}\iota\iota^{\T}\Bigr),
\end{equation}
where $I_J$ is the identity matrix and $\iota$ is the $J$-vector of ones. This matrix satisfies $\Sigma\iota = 0$, as adding up requires. It
turns maximum likelihood into least squares on the whole system. To see this, drop equation $J$. The
covariance matrix of the remaining errors is $\sigma^2(I_{J-1} - \frac{1}{J}\mathbf{1}\mathbf{1}^{\T})$,
where $\mathbf{1}$ is a $(J-1)$-vector of ones. Its inverse is $\sigma^{-2}(I_{J-1} +
\mathbf{1}\mathbf{1}^{\T})$, as multiplying the two matrices confirms:
\begin{equation*}
  \Bigl(I - \frac{1}{J}\mathbf{1}\mathbf{1}^{\T}\Bigr)\bigl(I + \mathbf{1}\mathbf{1}^{\T}\bigr)
  = I + \mathbf{1}\mathbf{1}^{\T} - \frac{1}{J}\mathbf{1}\mathbf{1}^{\T} - \frac{J-1}{J}\mathbf{1}\mathbf{1}^{\T}
  = I .
\end{equation*}
The quadratic form in the likelihood is then
\begin{equation*}
  \tilde\xi_t^{\T}\bigl(I + \mathbf{1}\mathbf{1}^{\T}\bigr)\tilde\xi_t
  = \sum_{i<J}\xi_{it}^2 + \Bigl(\sum_{i<J}\xi_{it}\Bigr)^2 = \sum_{i=1}^{J}\xi_{it}^2,
\end{equation*}
using $\xi_{Jt} = -\sum_{i<J}\xi_{it}$. Maximizing the likelihood over the coefficients therefore minimizes
the total sum of squared residuals across all $J$ equations, the trace of the residual cross-product
matrix rather than its determinant.

\subsection{The nonlinear AIDS}\label{sec:05-nl}

Substituting the translog index into \eqref{eq:05-stoch} gives the estimating equations of the exact model
\citep[eq.~15]{DeatonMuellbauer1980a}:
\begin{equation}\label{eq:05-nl-aids}
  w_{it} = (\alpha_i - \beta_i\alpha_0) + \sum_j\gamma_{ij}\ln p_{jt}
  + \beta_i\Bigl\{\ln I_t - \sum_k\alpha_k\ln p_{kt}
  - \frac{1}{2}\sum_k\sum_j\gamma_{kj}\ln p_{kt}\ln p_{jt}\Bigr\} + \xi_{it}.
\end{equation}
The system is nonlinear in the parameters because $\beta_i$ multiplies $\alpha_k$ and $\gamma_{kj}$. It can
be estimated by nonlinear least squares, by maximum likelihood, or by the generalized method of moments
(GMM) when instruments are needed, with or without the homogeneity and symmetry restrictions. The
computation is easier than it looks. Given $\beta$, the equations are linear in $(\alpha, \gamma)$.
Given $(\alpha, \gamma)$, the index is known and the equations are linear in $\beta$. The first-order
conditions can therefore be solved by alternating between two linear problems
\citep[p.~316]{DeatonMuellbauer1980a}.

All parameters of \eqref{eq:05-nl-aids} are identified given enough variation in prices and expenditure.
The constant $\alpha_0$ is the exception in practice. It is identified only through the appearance of the
$\alpha_k$ inside the braces. When individual prices move closely together, $\ln\mathbf{P}_t$ is
insensitive to its weights. A change in $\alpha_0$ can then be offset by changes in the $\alpha_i$ with
little effect on the fit. Deaton and Muellbauer recommend fixing $\alpha_0$ a priori. At unit prices
$\ln a(p) = \alpha_0$, so $e^{\alpha_0}$ is the cost of subsistence when all prices equal one. In their
words it is ``the outlay required for a minimal standard of living,'' and a plausible value is easy to
choose \citep[p.~316]{DeatonMuellbauer1980a}.

\subsection{The linear approximate AIDS and the Stone index}\label{sec:05-stone}

The same collinearity suggests a shortcut. If $\mathbf{P}_t$ is approximately proportional to a known
index $P_t^*$, say $\mathbf{P}_t \approx \varphi P_t^*$, then
\begin{equation}\label{eq:05-la-aids}
  w_{it} = (\alpha_i - \beta_i\ln\varphi) + \sum_j\gamma_{ij}\ln p_{jt}
  + \beta_i\ln\frac{I_t}{P_t^*} + \xi_{it}
\end{equation}
\citep[eq.~16]{DeatonMuellbauer1980a}. This \emph{linear approximate AIDS} (LA-AIDS) is linear in the
parameters. Without cross-equation restrictions it can be estimated by OLS equation by equation, which is
maximum likelihood for the system and satisfies adding up. The intercepts $\alpha_i^* = \alpha_i -
\beta_i\ln\varphi$ still sum to one because $\sum_i\beta_i = 0$. The usual choice of $P_t^*$ is the
\emph{Stone index}
\begin{equation}\label{eq:05-stone}
  \ln P_t^S = \sum_k w_{kt}\ln p_{kt}.
\end{equation}
It is exact, up to the constant $\alpha_0$, in one case. If $\gamma_{ij} = 0$ and $\beta_i = 0$ for all
$i, j$, then $w_k = \alpha_k$ and $\ln\mathbf{P} = \alpha_0 + \sum_k\alpha_k\ln p_k = \alpha_0 + \ln P^S$.
This is the Cobb-Douglas case. Otherwise the approximation is good when prices move together.

\casehead{The exact gap between the two models}
How good is the approximation, and what goes wrong when it fails? Substituting the share equations into
\eqref{eq:05-stone} and subtracting the translog index gives an exact expression for the gap:
\begin{align}
  \ln P_t^S - \ln\mathbf{P}_t
  &= \sum_k\Bigl(\alpha_k + \sum_j\gamma_{kj}\ln p_{jt} + \beta_k\ln\frac{I_t}{\mathbf{P}_t}
     + \xi_{kt}\Bigr)\ln p_{kt} - \ln\mathbf{P}_t \notag\\
  &= -\alpha_0 + \frac{1}{2}\sum_k\sum_j\gamma_{kj}\ln p_{kt}\ln p_{jt}
     + \ln\frac{I_t}{\mathbf{P}_t}\sum_k\beta_k\ln p_{kt} + \sum_k\xi_{kt}\ln p_{kt}.
  \label{eq:05-stone-gap}
\end{align}
The true model \eqref{eq:05-stoch} can be written in terms of the Stone index by adding and subtracting
$\beta_i\ln P^S_t$:
\begin{equation}\label{eq:05-la-exact}
\begin{split}
  w_{it} = {}& (\alpha_i - \beta_i\alpha_0) + \sum_j\gamma_{ij}\ln p_{jt} + \beta_i\ln\frac{I_t}{P^S_t} \\
  &+ \beta_i\Bigl[\frac{1}{2}\sum_k\sum_j\gamma_{kj}\ln p_{kt}\ln p_{jt}
  + \ln\frac{I_t}{\mathbf{P}_t}\sum_k\beta_k\ln p_{kt} + \sum_k\xi_{kt}\ln p_{kt}\Bigr] + \xi_{it}.
\end{split}
\end{equation}
The LA-AIDS omits the bracket. Each of its three terms explains one known problem of the Stone index.

\emph{Interpretation of the coefficients.} The middle term is linear in $\ln p_{kt}$ with coefficient
$\beta_i\beta_k\ln\frac{I_t}{\mathbf{P}_t}$. When real expenditure varies around a mean $m$, this term is
absorbed into the price coefficients, and the LA-AIDS estimates $\gamma_{ij}^{LA} \approx \gamma_{ij} +
\beta_i\beta_j m$ rather than $\gamma_{ij}$. The usual LA-AIDS elasticity formula
\begin{equation}\label{eq:05-la-elast}
  \varepsilon_{ij} \approx -\indic(i=j) + \frac{\gamma_{ij}^{LA} - \beta_i w_j}{w_i}
\end{equation}
then equals $-\indic(i=j) + \frac{\gamma_{ij} - \beta_i(w_j - \beta_j m)}{w_i}$, which is the exact AIDS
elasticity of \eqref{eq:05-elast-marshall} at $\ln\frac{I}{\mathbf{P}} = m$. Although the LA coefficients
differ from the AIDS coefficients, the LA elasticities can be close to the AIDS elasticities near the
center of the data.

\emph{Simultaneity.} The last term contains the errors of all equations. The Stone index is built from
current shares, so the regressor $\ln\frac{I_t}{P_t^S}$ is correlated with $\xi_{it}$. The term is small
when log prices are close to zero.

\emph{Units of measurement.} \citet{Moschini1995} shows that the Stone index is not invariant to the units
in which goods are measured. Suppose the unit of good $k$ is changed, so that $p_{kt}$ becomes
$\lambda p_{kt}$ ($\lambda > 1$ for a larger unit) and $\ln p_{kt}$ becomes $\ln p_{kt} + c$ with
$c = \ln\lambda$. The exact AIDS absorbs
the change. Replacing $\alpha_j$ by $\alpha_j - c\gamma_{jk}$ for every $j$ and $\alpha_0$ by $\alpha_0 -
c\alpha_k + \frac{1}{2}c^2\gamma_{kk}$ leaves both the translog index and every share equation unchanged
as functions of the new prices. To verify, expand the index in the new log prices $\ln p'_{jt}$:
\begin{equation*}
  \ln\mathbf{P}_t = \Bigl(\alpha_0 - c\alpha_k + \frac{1}{2}c^2\gamma_{kk}\Bigr)
  + \sum_j(\alpha_j - c\gamma_{jk})\ln p'_{jt} + \frac{1}{2}\sum_j\sum_l\gamma_{jl}\ln p'_{jt}\ln p'_{lt},
\end{equation*}
using symmetry, and note that $\gamma_{ik}\ln p_{kt} = \gamma_{ik}\ln p'_{kt} - c\gamma_{ik}$ in each share
equation. The new $\alpha_j$ still sum to one because $\sum_j\gamma_{jk} = 0$. Fitted shares and
elasticities do not change. The Stone index, by contrast, shifts by $c\,w_{kt}$, which varies over
observations. Redoing \eqref{eq:05-stone-gap} in the new units, the bracket in \eqref{eq:05-la-exact} gains
the term $c(w_{kt} - \alpha_k)$ plus a constant. The LA-AIDS therefore omits a time-varying term
$\beta_i c\, w_{kt}$ that is correlated with its regressors. Its estimates of $\beta_i$, $\gamma_{ij}$ and
the elasticities depend on the arbitrary choice of units, and the dependence grows with $|c|$.

\begin{example}[Units and the Stone index: a simulation]\label{ex:05-stone-mc}
We generate data from the exact AIDS with the running-example parameters and $\alpha_0 = 0$. Log prices
are $\ln p_{jt} = \pi_t + s_e e_{jt}$, with a common component $\pi_t \sim N(0, 0.3^2)$ and relative-price
noise $e_{jt} \sim N(0,1)$. Log expenditure is $\ln I_t = 1 + \pi_t + 0.3g_t$ with $g_t \sim N(0,1)$, and
share errors sum to zero, with standard deviations of about $0.004$. Design A has nearly collinear prices
($s_e = 0.02$) and design B has substantial relative-price variation ($s_e = 0.20$). Each sample has $T =
100$ observations, and we draw 500 samples. Four estimators impose adding up, homogeneity and symmetry,
and minimize the total sum of squared residuals. The first is the exact nonlinear AIDS and the second the
LA-AIDS with the Stone index. The third is the same LA-AIDS after good 1 is measured in units one hundred
times larger ($p_1$ multiplied by 100, $c = \ln 100 = 4.61$). The fourth is an LA-AIDS deflated by a fixed-weight index with
sample-mean shares. \Cref{tab:05-stone-mc}
reports the results, with elasticities evaluated at the base point of \cref{ex:05-running}.

Three lessons emerge. First, with prices normalized to one near the center of the data, the Stone LA-AIDS
recovers the elasticities in both designs. It does so even though the gap between the indices in design B
is of the order of 0.1 in logs. Its own-price coefficient $\hat\gamma_{11}^{LA}$ averages $-0.190$, while
the AIDS coefficient is $-0.200$. The difference is $\beta_1^2 m = 0.01$, exactly as
\eqref{eq:05-la-exact} predicts. Second, merely changing the units of
good 1 moves $\hat\beta_1$ from $-0.100$ to about $-0.069$ and the own-price elasticity of good 1 from
$-1.38$ to $-1.26$. Third, the fixed-weight index gives the same estimates in either unit, and its
elasticities match those of the exact model. In design A the average largest gap between the demeaned Stone and translog indices is
0.008, the same magnitude that Deaton and Muellbauer found in British data (\cref{sec:05-dm-application}).
\end{example}

\begin{table}[htbp]
\centering
\caption{Estimating the AIDS with exact and approximate price indices (simulation)}
\label{tab:05-stone-mc}
\small
\begin{tabular}{lcccc}
\toprule
Estimator & $\hat\beta_1$ & $\hat\gamma_{11}$ & $\hat\varepsilon_{11}$ & $\hat\varepsilon_{33}$ \\
\midrule
True values & $-0.100$ & $-0.200$ & $-1.375$ & $-1.738$ \\
\addlinespace
\multicolumn{5}{l}{\textit{Design A: nearly collinear prices ($s_e = 0.02$)}} \\
Exact AIDS (nonlinear) & $-0.100$ (0.001) & $-0.200$ (0.017) & $-1.376$ (0.042) & $-1.739$ (0.066) \\
LA-AIDS, Stone index & $-0.100$ (0.001) & $-0.190$ (0.017) & $-1.376$ (0.042) & $-1.740$ (0.066) \\
LA-AIDS, Stone, rescaled $p_1$ & $-0.069$ (0.001) & $-0.130$ (0.012) & $-1.257$ (0.029) & $-1.667$ (0.060) \\
LA-AIDS, fixed-weight index & $-0.100$ (0.001) & $-0.190$ (0.017) & $-1.376$ (0.042) & $-1.739$ (0.066) \\
\addlinespace
\multicolumn{5}{l}{\textit{Design B: independent relative prices ($s_e = 0.20$)}} \\
Exact AIDS (nonlinear) & $-0.100$ (0.001) & $-0.200$ (0.002) & $-1.375$ (0.004) & $-1.738$ (0.007) \\
LA-AIDS, Stone index & $-0.100$ (0.001) & $-0.190$ (0.002) & $-1.375$ (0.005) & $-1.737$ (0.007) \\
LA-AIDS, Stone, rescaled $p_1$ & $-0.068$ (0.001) & $-0.130$ (0.001) & $-1.257$ (0.004) & $-1.665$ (0.007) \\
LA-AIDS, fixed-weight index & $-0.100$ (0.001) & $-0.190$ (0.002) & $-1.375$ (0.005) & $-1.738$ (0.007) \\
\bottomrule
\end{tabular}

\smallskip
\begin{minipage}{0.94\textwidth}
\footnotesize Notes: means over 500 simulated samples of $T = 100$, standard deviations across samples in
parentheses. Data are generated from the exact AIDS with the parameters of \cref{ex:05-running}.
Elasticities are evaluated at $p = (1,1,1)$, $\ln I = 1$, with the true shares $(0.40, 0.34, 0.26)$. The
LA-AIDS elasticities use \eqref{eq:05-la-elast}. The fixed-weight estimates are identical in either unit.
Computed by \texttt{figures/ch05/src/stone\_index\_mc.py}.
\end{minipage}
\end{table}

\casehead{Elasticities from the LA-AIDS}
Applied work often computes LA-AIDS elasticities from \eqref{eq:05-la-elast}, which treats
$\frac{\partial\ln P^S}{\partial\ln p_j}$ as $w_j$. \citet{GreenAlston1990} point out that the Stone
weights themselves respond to prices. Within the LA model,
\begin{equation*}
  \frac{\partial\ln P^S}{\partial\ln p_j} = w_j + \sum_k\ln p_k\frac{\partial w_k}{\partial\ln p_j}.
\end{equation*}
Let $\mathcal{D}$ be the matrix of share derivatives $\frac{\partial w_i}{\partial\ln p_j}$, $\ell$ the vector of
log prices, $\beta$ the vector of expenditure coefficients and $w$ the vector of shares. Differentiating
\eqref{eq:05-la-aids},
\begin{equation*}
  \mathcal{D} = \Gamma^{LA} - \beta w^{\T} - \beta\ell^{\T}\mathcal{D}
  \quad\Longrightarrow\quad
  \mathcal{D} = \bigl(I + \beta\ell^{\T}\bigr)^{-1}\bigl(\Gamma^{LA} - \beta w^{\T}\bigr).
\end{equation*}
By the Sherman-Morrison formula, $(I + \beta\ell^{\T})^{-1} = I - \frac{\beta\ell^{\T}}{1 +
\ell^{\T}\beta}$. At a point where all prices equal one, $\ell = 0$ and \eqref{eq:05-la-elast} is exact for
the LA model. Away from it the correction matters, and it depends on units for the reason just given.
Normalizing prices to one at the point of evaluation removes the problem.

\casehead{Remedies}
Several fixes are available. A \emph{Laspeyres-type} index with fixed weights, $\ln P_t^L = \sum_k\bar
w_k\ln p_{kt}$ (base-period or sample-mean shares), contains no current shares, so it removes the
simultaneity. A change of units shifts it by a constant that the intercepts absorb. The
\emph{Tornqvist} index $\ln P_t^T = \sum_k\frac{1}{2}(w_{kt} + w_k^0)\ln\frac{p_{kt}}{p_k^0}$ and the
log-linear analogue of the Paasche index, $\sum_k w_{kt}\ln\frac{p_{kt}}{p_k^0}$, use prices relative to a
base point $p^0$ and are therefore invariant to units \citep{Moschini1995}.

The Tornqvist index also has a basis in demand theory \citep[\S 7.1, eqs.~1.8--1.10,
pp.~174--175]{DeatonMuellbauer1980b}. Along any price path at fixed utility, Shephard's lemma gives
$\dd\ln e(p, u) = \sum_k w_k(p, u)\,\dd\ln p_k$, with Hicksian shares $w_k(p, u)$. Integrating from $p^0$
to $p^1$,
\begin{equation*}
  \ln\frac{e(p^1, u)}{e(p^0, u)} = \int_{p^0}^{p^1}\sum_k w_k(p, u)\,\dd\ln p_k ,
\end{equation*}
which is the log of the exact cost-of-living index at $u$. The \emph{Divisia index} replaces the Hicksian
shares in this integral by observed shares. Unless preferences are homothetic, observed shares move with
utility along the path, so the Divisia index depends on the path. Prices are observed only at discrete
dates, so the integral must be approximated. The Tornqvist index is a natural approximation, since it
averages the shares at the two ends. Following \citet{Diewert1976}, it equals the exact index at a
reference utility $u^*$ when the log expenditure function is quadratic in log prices and log utility. The
reference is then $u^* = \sqrt{u^0u^1}$, the geometric mean of the utilities at the two dates \citep[p.~174]{DeatonMuellbauer1980b}.
The reason is an identity for quadratic functions. Their change between two points equals the average of
the end-point gradients times the change in the arguments. At a fixed reference utility, the price
gradients are the Hicksian shares at that utility. They differ from the observed shares by terms linear
in the utility gap, which cancel when $\ln u^*$ is the midpoint of $\ln u^0$ and $\ln u^1$. If $\ln e$ is
quadratic in log prices and in $u$ itself, the reference is the arithmetic mean $\frac{1}{2}(u^0 + u^1)$
instead. A fixed-weight index shares the units invariance but not this link to the
exact index. A Stone index with lagged
shares removes the contemporaneous correlation with $\xi_t$ when the errors are serially uncorrelated.
Finally, we can iterate toward the exact model, as \citet[p.~321]{DeatonMuellbauer1980a} suggest:
\begin{enumerate}[label=\arabic*.]
  \item Compute $P^*_t$ (for example the Stone index) and estimate \eqref{eq:05-la-aids}.
  \item Compute $\ln\mathbf{P}_t$ from the estimated $\alpha$ and $\gamma$ and a chosen $\alpha_0$.
  \item Re-estimate the share equations with $\ln\mathbf{P}_t$ in place of $\ln P^*_t$.
  \item Repeat steps 2 and 3 until the estimates stop changing.
\end{enumerate}
Deaton and Muellbauer describe this as ``a computationally efficient way of obtaining good estimates of
the full nonlinear system.'' They also stress that \eqref{eq:05-la-aids} exists only as an approximation.
Once the restrictions have been chosen, the full system \eqref{eq:05-nl-aids} should be re-estimated
\citep[p.~317]{DeatonMuellbauer1980a}.

\subsection{First differences and the Rotterdam model}\label{sec:05-rotterdam}

Time-series applications often estimate in first differences. Differencing \eqref{eq:05-stoch} gives
\begin{equation}\label{eq:05-fd-aids}
  \Delta w_{it} = \beta_i\Delta\ln\frac{I_t}{\mathbf{P}_t} + \sum_j\gamma_{ij}\Delta\ln p_{jt}
  + \Delta\xi_{it}
\end{equation}
\citep[eq.~17]{DeatonMuellbauer1980a}, in which the $\alpha_i$ appear only through $\Delta\ln\mathbf{P}_t$.
Written out, $\Delta\ln\mathbf{P}_t$ involves $\sum_k\alpha_k\Delta\ln p_{kt}$ and $\frac{1}{2}\sum_k\sum_j
\gamma_{kj}\Delta(\ln p_{kt}\ln p_{jt})$ (their eq.~18). With nearly collinear prices the $\gamma_{ij}$ and
the products $\beta_i\alpha_j$ are hard to separate. In practice $\Delta\ln\mathbf{P}_t$ is replaced by
the change in a Stone index or by $\sum_k w_{kt}\Delta\ln p_{kt}$.

With the second choice the right-hand side becomes identical to that of the \emph{Rotterdam model} of
\citet{Theil1965} and \citet{Barten1969}. The Rotterdam model is a differential demand system. Totally
differentiating $q_i(p, I)$ and multiplying by $w_i$ gives $w_i\,\dd\ln q_i = w_i\eta_i\,\dd\ln I +
\sum_j w_i\varepsilon_{ij}\,\dd\ln p_j$. Substituting the Slutsky equation in elasticity form,
$\varepsilon_{ij} = \varepsilon^h_{ij} - \eta_i w_j$ (\cref{sec:05-elasticity-identities}), gives
\begin{equation}\label{eq:05-rotterdam}
  w_{it}\Delta\ln q_{it} = b_i\Bigl(\Delta\ln I_t - \sum_k w_{kt}\Delta\ln p_{kt}\Bigr)
  + \sum_j c_{ij}\Delta\ln p_{jt},
\end{equation}
with $b_i = w_i\eta_i = p_i\frac{\partial q_i}{\partial I}$, the marginal budget share, and $c_{ij} =
w_i\varepsilon^h_{ij} = \frac{p_ip_jS_{ij}}{I}$. The Rotterdam model treats $b_i$ and $c_{ij}$ as constant
parameters \citep[eq.~19]{DeatonMuellbauer1980a}.

The two models differ in the dependent variable. The AIDS has $\Delta w_i$, which is approximately $w_i
\Delta\ln w_i = w_i(\Delta\ln q_i + \Delta\ln p_i - \Delta\ln I)$. So the first-difference AIDS replaces the
Rotterdam dependent variable $w_i\Delta\ln q_i$ by $w_i\Delta\ln w_i$, adding $w_i\Delta\ln\frac{p_i}{I}$.
Both models are linear in constant parameters, and both test homogeneity and symmetry through linear
restrictions. Their parameters mean different things. The Rotterdam $c_{ij}$ are compensated price effects,
so negativity applies directly to the matrix $[c_{ij}]$. In the AIDS negativity applies to the matrix $K$
of \eqref{eq:05-K}, which depends on the data. Above all, the AIDS is derived from explicit preferences.
For forecasting demand this difference may matter little. For cost-of-living indices, equivalence scales
or welfare calculations, the link between estimated parameters and preferences matters a great deal
\citep[pp.~317--318]{DeatonMuellbauer1980a}.

\subsection{Testing the restrictions of demand theory}\label{sec:05-testing}

The restrictions of \cref{prop:05-restrictions} have different statistical status.

\emph{Adding up} cannot be tested with share data. The shares sum to one by construction, and OLS
estimates satisfy adding up automatically, as shown in \cref{sec:05-singular}.

\emph{Homogeneity}, $\sum_j\gamma_{ij} = 0$, involves one equation at a time and can be tested equation by
equation. Because the Stone index is homogeneous of degree one in prices, homogeneity can be imposed on
\eqref{eq:05-la-aids} by writing prices relative to the price of good $J$:
\begin{equation}\label{eq:05-homog-sub}
  w_{it} = \alpha_i^* + \sum_{j=1}^{J-1}\gamma_{ij}\ln\frac{p_{jt}}{p_{Jt}} + \beta_i\ln\frac{I_t}{P^*_t}
  + \xi_{it}
\end{equation}
\citep[eq.~20]{DeatonMuellbauer1980a}. The identity $\sum_j\gamma_{ij}\ln p_j = \sum_{j<J}\gamma_{ij}
\ln\frac{p_j}{p_J} + (\sum_j\gamma_{ij})\ln p_J$ shows that \eqref{eq:05-homog-sub} is the unrestricted
equation with the coefficient on $\ln p_J$ set to $\sum_j\gamma_{ij} = 0$. An $F$ test compares the
residual sums of squares of the two regressions. Equivalently, it is the square of the $t$ statistic on
$\sum_j\hat\gamma_{ij}$ in the unrestricted regression.

\emph{Symmetry}, $\gamma_{ij} = \gamma_{ji}$, links equations and can only be tested on the system, for
example by a likelihood-ratio test,
\begin{equation*}
  LR = 2\bigl(\ln L_{\mathrm{U}} - \ln L_{\mathrm{R}}\bigr) \xrightarrow{d} \chi^2_r .
\end{equation*}
Here $L_{\mathrm{U}}$ and $L_{\mathrm{R}}$ are the maximized likelihoods without and with the
restrictions, and $r$ is the number of restrictions. In a $J$-good system estimated on $J - 1$ equations, homogeneity
imposes $J - 1$ restrictions and symmetry imposes $\frac{(J-1)(J-2)}{2}$ more. For $J = 8$ these numbers
are 7 and 21. The test is asymptotic, and allowance must be made for that when samples are short.

Deaton and Muellbauer test symmetry whether or not homogeneity has been rejected. They report an
objection by Grayham Mizon, who argued that optimal inference stops testing once a rejection occurs.
Their reply is that the choice of maintained hypothesis is partly a matter of judgment. Many economists
would treat the absence of money illusion as maintained and regard symmetry as the interesting test.
Moreover, ``few if any tests in econometrics are carried out within the framework of maintained hypotheses which
are even widely accepted'' \citep[pp.~318--319]{DeatonMuellbauer1980a}.

\emph{Negativity} cannot be imposed by linear restrictions and is checked after estimation. At each
observation we compute the matrix $K$ of \eqref{eq:05-K} and its eigenvalues. One eigenvalue is
identically zero because $K\iota = 0$, and the others should be nonpositive. Negativity can be imposed
locally, at one point, by a Cholesky reparameterization of the Slutsky matrix \citep{Lau1978,Moschini1999}.

\subsection{Case study: the AIDS on British data, 1954--1974}\label{sec:05-dm-application}

\casehead{Question}
The empirical half of \citet{DeatonMuellbauer1980a} asks whether aggregate consumer demand obeys the
homogeneity, symmetry and negativity restrictions of demand theory. Earlier studies with the Rotterdam and
translog models had rejected homogeneity. The AIDS lets the question be posed within a maintained model
flexible enough to approximate any demand system to first order (\cref{sec:05-flexible}). A rejection
then speaks against the theory or its static formulation, and an arbitrary functional form cannot take
the blame. A second question is practical. Does the linear approximation with the Stone
index work well enough to serve as a routine tool? \Cref{tab:05-dm-map} at the end of this case study maps
its data, identifying variation and maintained assumptions.

\casehead{Data}
The data are 21 annual British observations, from 1954 to 1974 inclusive. They cover eight groups of
nondurable consumers' expenditure: food, clothing, housing services, fuel, drink and tobacco, transport
and communication services, other goods, and other services. For each year and group the data give a
budget share and a price index, together with total expenditure on the eight groups. Over the period the
price level rose by a factor of about 2.6. In their Table 3 the exact index rises from 0.571 in 1954 to
1.461 in 1974, with 1970 equal to one.

\casehead{Model}
The estimating model is the aggregate AIDS \eqref{eq:05-aids-aggregate} of a representative household. Deaton
and Muellbauer assume that the aggregation index $\kappa$ is constant, or that its deviations are
independent of mean expenditure and prices. The terms $-\beta_i\ln\kappa$ are then absorbed in the
intercepts, as in \cref{sec:05-almost-exact}.

\casehead{Estimating equations and estimator}
The procedure has five steps.
\begin{enumerate}[label=\arabic*.]
  \item Estimate the LA-AIDS \eqref{eq:05-la-aids} with the Stone index, equation by equation by OLS,
  with no restrictions beyond adding up, which holds automatically (their Table 1).
  \item Impose homogeneity by substitution, \eqref{eq:05-homog-sub}, re-estimate by OLS, and compute an
  $F$ test for each equation.
  \item Impose symmetry, replace $P^*$ by the exact index $\mathbf{P}$ with $\alpha_0$ fixed a priori, and
  estimate the nonlinear system \eqref{eq:05-nl-aids} by maximum likelihood under the covariance structure
  \eqref{eq:05-dm-cov}. Under symmetry the system has fourteen free intercepts and expenditure
  coefficients and twenty-eight free price coefficients, the upper triangle, diagonal included, of the
  $\gamma$ matrix with its last row and column deleted.
  \item Compare $P^*$ with $\mathbf{P}$, re-estimate the unrestricted and homogeneous models with
  $\mathbf{P}$, and check concavity with the eigenvalues of $K$ in every year.
  \item Repeat the analysis in first differences, adding intercepts $d_i$ to \eqref{eq:05-fd-aids}:
  \begin{equation}\label{eq:05-dm-fd}
    \Delta w_{it} = d_i + \beta_i\Delta\ln\frac{I_t}{\mathbf{P}_t} + \sum_k\gamma_{ik}\Delta\ln p_{kt}
    + \Delta\xi_{it}
  \end{equation}
  (their eq.~21, where the intercepts are written $\eta_i$). Collinearity prevented linking $\mathbf{P}$
  to the first-difference estimates, so the index computed from the symmetric levels estimates is used
  throughout.
\end{enumerate}

\casehead{Identification}
Prices and total expenditure are treated as exogenous. The article does not discuss this assumption,
which is common in the analysis of national aggregates. The authors' book does, and doubts both halves
of it \citep[\S 3.5, pp.~81--82]{DeatonMuellbauer1980b}. On total expenditure it says that ``individual
expenditures are as much a determinant of the total as the other way round'' (p.~81). On prices it says
that demand as a function of prices ``only makes sense if prices are set exogenously (say by
manufacturers) and quantities supplied elastically to meet whatever demand emerges.'' That is probably
true of most goods, ``although food may be a partial exception'' (pp.~81--82). The book links this to
the difficulty of estimating a ``sensible'' price elasticity for food. The positive own-price elasticity of
food in the unconstrained levels model of \cref{tab:05-dm-table2} (0.07) fits this warning. The case for
treating the regressors as given rests on two premises. Broad group prices are driven largely by costs and
by the general price level. Total nondurable expenditure is taken as given at the second stage of an
intertemporal budgeting problem. That second premise is the weak intertemporal separability that the
authors question in their interpretation. The price coefficients are identified by movements in
relative prices, and the expenditure coefficients by movements in real expenditure. Relative prices moved
little, which explains both the closeness of $P^*$ and $\mathbf{P}$ and the imprecision of many
$\gamma_{ij}$.

\casehead{Unconstrained estimates and homogeneity}
\Cref{tab:05-dm-table1} reproduces selected columns of their Table 1. The expenditure coefficients
classify food ($\beta = -0.160$) and housing ($-0.104$) as necessities and the other six groups as
luxuries. Twenty-two of the sixty-four price coefficients have $t$ values above two in absolute value.
Yet no variable has a detectable effect on the share of fuel, and few affect the shares of other goods
and other services. The column $\sum_j\gamma_{ij}$ gives 100 times the effect on each share of a 1
percent rise in all prices and total expenditure. Homogeneity requires it to be zero, and its $t$ value
is the square root of the $F$ statistic for the restriction. The authors' book reports these $F$
statistics, with one and eleven degrees of freedom, together with the fit of the free and the
homogeneous equations \citep[\S 3.4, Table~3.4, pp.~77--78]{DeatonMuellbauer1980b}.
\Cref{tab:05-dm-homog} reproduces that table. The rejection thresholds are 4.84 at the 5 percent level and
9.65 at the 1 percent level (p.~77). Homogeneity is rejected at 1 percent for food ($F = 19.4$, $t =
-4.4$), clothing ($20.3$, $-4.5$), housing ($82.8$, $9.1$) and transport and communication ($171.6$,
$13.1$). A proportional increase in all prices and expenditure would lower the budget shares of food and
clothing and raise those of housing and of transport and communication.

The fit is very high in both versions, with $R^2$ above 0.95 in five of the eight free equations and
never below 0.84. Budget
shares are harder to predict than quantities or expenditures, so these values are informative. For
comparison, the linear expenditure system fitted to the same share data gives $R^2 = -0.16$ for other
services (p.~77). A negative $R^2$ is possible because the restrictions of that system fit the share worse
than its sample mean. The \emph{Durbin-Watson statistic} tests for first-order serial correlation in the
residuals. Values near 2 indicate none, and values well below 2 indicate positive correlation. In each of
the four equations that reject homogeneity, imposing it makes the statistic fall sharply, for example
from 2.24 to 1.36 for transport and communication. For the other four groups the statistic barely moves.
The imposition of homogeneity thus generates positive serial correlation. Deaton and Muellbauer read this
as a sign that the rejection may come from misspecified dynamics (p.~77).

\begin{table}[htbp]
\centering
\caption{Unconstrained LA-AIDS estimates and tests of homogeneity, British data, 1954--1974}
\label{tab:05-dm-table1}
\small
\begin{tabular}{lcccc}
\toprule
Group & $\alpha_i^*$ & $\beta_i$ & $\gamma_{ii}$ & $\sum_j\gamma_{ij}$ \\
\midrule
Food & 1.221 & $-0.160$ & 0.186 & $-0.033$ \\
 & (7.4) & ($-6.1$) & (9.8) & ($-4.4$) \\
Clothing & $-0.482$ & 0.091 & 0.016 & $-0.032$ \\
 & ($-3.1$) & (3.7) & (1.0) & ($-4.5$) \\
Housing & 0.793 & $-0.104$ & 0.088 & 0.051 \\
 & (6.3) & ($-5.1$) & (7.2) & (9.1) \\
Fuel & $-0.159$ & 0.033 & 0.037 & $-0.010$ \\
 & ($-0.8$) & (1.0) & (1.6) & ($-1.1$) \\
Drink and tobacco & $-0.043$ & 0.028 & 0.056 & 0.001 \\
 & ($-0.3$) & (1.2) & (6.9) & (0.0) \\
Transport and communication & $-0.061$ & 0.029 & $-0.023$ & 0.040 \\
 & ($-0.9$) & (2.6) & ($-2.2$) & (13.1) \\
Other goods & $-0.038$ & 0.022 & 0.032 & $-0.005$ \\
 & ($-0.2$) & (0.9) & (1.5) & (0.7) \\
Other services & $-0.231$ & 0.060 & 0.019 & $-0.014$ \\
 & ($-1.5$) & (2.4) & (0.9) & ($-2.0$) \\
\bottomrule
\end{tabular}

\smallskip
\begin{minipage}{0.94\textwidth}
\footnotesize Notes: selected columns of \citet[Table~1]{DeatonMuellbauer1980a}, with $t$ values in
parentheses. OLS estimates of \eqref{eq:05-la-aids} with the Stone index, equation by equation. The column
$\gamma_{ii}$ is the own-price coefficient. The original also reports the seven cross-price coefficients
of each equation, equation standard errors, $R^2$ and Durbin-Watson statistics, which
\cref{tab:05-dm-homog} gives. $\sum_j\gamma_{ij}$ is the row sum of the unconstrained price
coefficients, and its $t$ value tests homogeneity. The $t$ value for other goods is printed as positive in
the original. For drink and tobacco the printed $t$ of 0.0 does not match the book's $F$ of 0.2 in
\cref{tab:05-dm-homog}, which implies $|t| \approx 0.45$. Both are far from significance. Each equation
uses $T = 21$ annual observations.
\end{minipage}
\end{table}

\begin{table}[htbp]
\centering
\caption{Tests of homogeneity on British data, 1954--1974}
\label{tab:05-dm-homog}
\small
\begin{tabular}{lccccc}
\toprule
 & \multicolumn{2}{c}{Unrestricted} & \multicolumn{2}{c}{Homogeneous} & Test \\
\cmidrule(lr){2-3}\cmidrule(lr){4-5}\cmidrule(lr){6-6}
Group & $R^2$ & Durbin-Watson & $R^2$ & Durbin-Watson & $F(1, 11)$ \\
\midrule
Food & 0.999 & 2.33 & 0.998 & 1.74 & 19.4 \\
Clothing & 0.984 & 2.29 & 0.955 & 1.55 & 20.3 \\
Housing & 0.999 & 1.89 & 0.992 & 1.29 & 82.8 \\
Fuel & 0.883 & 2.25 & 0.870 & 2.03 & 1.2 \\
Drink and tobacco & 0.969 & 2.96 & 0.969 & 2.93 & 0.2 \\
Transport and communication & 1.000 & 2.24 & 0.992 & 1.36 & 171.6 \\
Other goods & 0.885 & 1.92 & 0.880 & 1.91 & 0.5 \\
Other services & 0.843 & 2.27 & 0.788 & 1.98 & 4.0 \\
\bottomrule
\end{tabular}

\smallskip
\begin{minipage}{0.94\textwidth}
\footnotesize Notes: reproduced from \citet[Table~3.4, p.~78]{DeatonMuellbauer1980b}. Equation-by-equation
OLS estimates of the LA-AIDS with the Stone index, $T = 21$, the same regressions as
\cref{tab:05-dm-table1}. $F(1, 11)$ tests the single homogeneity restriction of each equation. The
thresholds are 4.84 at 5 percent and 9.65 at 1 percent. Except for drink and tobacco, each $F$ equals the
square of the $t$ value in \cref{tab:05-dm-table1} up to rounding.
\end{minipage}
\end{table}

\Cref{tab:05-dm-table2} reproduces their Table 2, the total expenditure and own-price elasticities of
all four specifications. In the unconstrained levels model, food has a positive own-price elasticity
(0.07), which is insignificant. Demand is generally price inelastic, and only transport and
communication is price elastic ($-1.21$). The elasticities that change most when homogeneity is imposed
belong to the four groups that reject it. The expenditure elasticity of housing rises from 0.30 to 0.79,
and that of transport and communication from 1.23 to 1.73.

\begin{table}[htbp]
\centering
\caption{Total expenditure and own-price elasticities, British data, 1954--1974}
\label{tab:05-dm-table2}
\small
\begin{tabular}{lcccccccc}
\toprule
 & \multicolumn{4}{c}{Levels model} & \multicolumn{4}{c}{First-differences model} \\
\cmidrule(lr){2-5}\cmidrule(lr){6-9}
 & \multicolumn{2}{c}{Unconstrained} & \multicolumn{2}{c}{Homogeneous}
 & \multicolumn{2}{c}{Unconstrained} & \multicolumn{2}{c}{Homogeneous} \\
\cmidrule(lr){2-3}\cmidrule(lr){4-5}\cmidrule(lr){6-7}\cmidrule(lr){8-9}
Group & $\eta_i$ & $\varepsilon_{ii}$ & $\eta_i$ & $\varepsilon_{ii}$ & $\eta_i$ & $\varepsilon_{ii}$
 & $\eta_i$ & $\varepsilon_{ii}$ \\
\midrule
Food & 0.21 & 0.07 & 0.04 & $-0.01$ & 0.04 & 0.22 & 0.17 & $-0.00$ \\
Clothing & 2.00 & $-0.92$ & 1.51 & $-0.48$ & 2.83 & $-0.94$ & 2.92 & $-0.94$ \\
Housing & 0.30 & $-0.31$ & 0.79 & $-0.16$ & 0.04 & $-0.31$ & $-0.02$ & $-0.30$ \\
Fuel & 1.67 & $-0.28$ & 1.37 & 0.10 & 1.00 & 0.00 & 0.86 & $-0.08$ \\
Drink and tobacco & 1.22 & $-0.60$ & 1.22 & $-0.62$ & 1.37 & $-0.67$ & 1.36 & $-0.68$ \\
Transport and comm. & 1.23 & $-1.21$ & 1.73 & $-0.92$ & 1.14 & $-1.23$ & 1.05 & $-1.17$ \\
Other goods & 1.21 & $-0.72$ & 1.15 & $-0.77$ & 2.03 & $-0.52$ & 1.92 & $-0.47$ \\
Other services & 1.40 & $-0.93$ & 1.28 & $-0.78$ & 1.03 & $-0.78$ & 1.06 & $-0.74$ \\
\bottomrule
\end{tabular}

\smallskip
\begin{minipage}{0.94\textwidth}
\footnotesize Notes: reproduced from \citet[Table~2]{DeatonMuellbauer1980a}, who write $e_i$ and $e_{ii}$.
$\eta_i$ is the total expenditure elasticity and $\varepsilon_{ii}$ the uncompensated own-price
elasticity.
\end{minipage}
\end{table}

\casehead{The Stone index versus the exact index}
Before imposing symmetry with the exact index, Deaton and Muellbauer compare the two indices year by year
(their Table 3, both normalized to one in 1970). The largest absolute difference is 0.008, reached in 1955
(0.587 against 0.595) and 1973 (1.271 against 1.279). At the ends of the sample the indices are 0.566 and
0.571 in 1954 and 1.465 and 1.461 in 1974. Re-estimating the unrestricted and homogeneous models with
$\mathbf{P}$ in place of $P^*$ changes the likelihoods little compared with the effect of the
restrictions, as the paired columns of \cref{tab:05-dm-table4} show. In the authors' words, it
``confirmed the empirical unimportance of the difference.'' They add that this finding is conditional on the kind of relative
price movements in their sample.

\casehead{Symmetry}
\Cref{tab:05-dm-table4} reproduces their Table 4 and computes the likelihood-ratio statistics. In levels
with the exact index, homogeneity gives $LR = 1723.8 - 1585.1 = 138.7$ on 7 degrees of freedom, far
above the 5 percent threshold of 14.07. Symmetry given homogeneity gives $LR = 1585.1 - 1491.0 =
94.1$ on 21 degrees of freedom, against a 5 percent threshold of 32.67. Symmetry is therefore rejected whether
or not homogeneity is part of the maintained hypothesis. The authors note that the meaning of this
rejection is unclear without an explanation of the failure of homogeneity. For example, habits that are
ignored can destroy symmetry even when preferences satisfy it. Their book reads the same evidence more
softly. Imposing symmetry led to ``a further but less sharp drop in likelihood,'' especially in first
differences \citep[\S 3.4, p.~78]{DeatonMuellbauer1980b}. Hence ``it is unclear whether we must reject
symmetry as a restriction beyond homogeneity'' (p.~78). Looking across models and countries, the book concludes
that ``demand functions fitted to aggregate time series data are not homogeneous and probably not
symmetric'' (p.~78).

\begin{table}[htbp]
\centering
\caption{Likelihood values and likelihood-ratio tests, British data, 1954--1974}
\label{tab:05-dm-table4}
\small
\begin{tabular}{lcccc}
\toprule
\multicolumn{5}{l}{\textit{Panel A: values of $2\ln L$}} \\
 & \multicolumn{2}{c}{Levels} & \multicolumn{2}{c}{First differences} \\
\cmidrule(lr){2-3}\cmidrule(lr){4-5}
 & Using $P^*$ & Using $\mathbf{P}$ & Using $P^*$ & Using $\mathbf{P}$ \\
\midrule
Unrestricted & 1722.5 & 1723.8 & 1560.0 & 1560.3 \\
Homogeneous & 1579.6 & 1585.1 (7) & 1546.6 & 1547.9 (7) \\
Symmetric & & 1491.0 (21) & & 1508.8 (21) \\
\midrule
\multicolumn{5}{l}{\textit{Panel B: likelihood-ratio tests computed from Panel A, using $\mathbf{P}$}} \\
 & $LR$ & df & 5\% threshold & $p$-value \\
\midrule
\multicolumn{5}{l}{Levels} \\
\quad Homogeneity & 138.7 & 7 & 14.07 & $< 10^{-20}$ \\
\quad Symmetry given homogeneity & 94.1 & 21 & 32.67 & $3 \times 10^{-11}$ \\
\multicolumn{5}{l}{First differences} \\
\quad Homogeneity & 12.4 & 7 & 14.07 & 0.088 \\
\quad Symmetry given homogeneity & 39.1 & 21 & 32.67 & 0.010 \\
\bottomrule
\end{tabular}

\smallskip
\begin{minipage}{0.94\textwidth}
\footnotesize Notes: Panel A reproduces \citet[Table~4]{DeatonMuellbauer1980a}, with numbers of restrictions in
parentheses. Values can be compared within columns only. Panel B: $LR$ is the difference in $2\ln L$
between the less and the more restricted model, compared with the $\chi^2$ distribution with df degrees of
freedom. Computed by \texttt{figures/ch05/src/dm\_lr\_tests.py}.
\end{minipage}
\end{table}

\casehead{Negativity}
The symmetric estimates themselves are not reported. Deaton and Muellbauer evaluate the matrix $K$ of
\eqref{eq:05-K} in every year. One eigenvalue is identically zero, and concavity requires the others to be
negative. They find one positive eigenvalue early in the sample and two by its end. The clearest symptom
is a positive compensated own-price elasticity for fuel throughout the period. They argue that the
failure matters even though the symmetric model has already been rejected. Price and quantity index
numbers and optimal tax formulas require cost functions that are homogeneous and concave. Applied
welfare work had therefore relied on the linear expenditure system, whose concavity is virtually
guaranteed when no good is inferior, but whose restrictiveness is known to be empirically false.

\casehead{Robustness: first differences}
The first-difference model \eqref{eq:05-dm-fd} changes the picture. Homogeneity is now rejected only for
food and for transport and communication. For clothing and housing the $F$ ratios become insignificant,
and in both equations the intercept $d_i$, insignificant without constraints, becomes significant when
homogeneity is imposed. For transport and communication the intercept also becomes significant, but the
$F$ ratio stays significant. An intercept in first differences is a time trend in levels, so these
results point to trends, stocks or other omitted variables. The first-difference expenditure
elasticities are higher than the levels estimates where stock effects are likely (clothing, other goods)
and lower where short-run expenditure effects should be limited (food, housing, transport). For the
system, homogeneity is no longer rejected at 5 percent ($LR = 12.4$ against 14.07). Symmetry given
homogeneity is ``only just rejected'' ($LR = 39.1$ against 32.67, $p = 0.010$). Concavity tests still
reveal several violations.

\casehead{Interpretation}
Deaton and Muellbauer read the rejection of homogeneity, together with the serial correlation it
induces, as a sign of dynamic misspecification. They offer several explanations and say that the data
cannot discriminate among them.
\begin{enumerate}[label=(\roman*)]
  \item Spending on some items is inflexible in the short run, housing being the obvious case. Stocks,
  lagged dependent variables or time trends, which the absolute price level may proxy, are omitted.
  \item Price expectations are omitted. The frequency with which goods are bought should affect how
  spending responds to prices when relative or absolute prices change rapidly.
  \item The aggregation index $\kappa$, which reflects the distribution of budgets and the demographic
  structure, may be correlated with mean expenditure and prices (\cref{sec:05-almost-exact}).
  \item Weak intertemporal separability of nondurables, which the static model requires, may fail.
\end{enumerate}
They conclude that the AIDS explains a high proportion of the variance of budget shares. Without
arbitrary time trends, however, it does so ``in a way which is inconsistent with the hypothesis of
consumers making decisions according to the model's demand functions''
\citep[pp.~322--323]{DeatonMuellbauer1980a}. With trends, much of the conflict between data
and theory disappears.

\casehead{Limitations and critiques}
The application has the weaknesses of its time and data. Twenty-one observations support a system with
dozens of parameters, so the covariance matrix has to be restricted to \eqref{eq:05-dm-cov}, and the
likelihood-ratio tests rest on asymptotic approximations in a short sample. The article treats prices and
expenditure as exogenous without discussion, although the authors' book doubts both
\citep[\S 3.5, pp.~81--82]{DeatonMuellbauer1980b}. The aggregation index is assumed constant, although the authors
themselves list its variation as a possible cause of the rejections. The symmetric estimates violate
concavity, and the authors leave the construction of a concave version to future work. Later research
took up each point. It added demographic and distributional variables to aggregate demand and built
dynamic and habit-forming versions of the AIDS. It also developed richer Engel curves, such as the
quadratic AIDS of \citet{BanksBlundellLewbel1997} (\cref{sec:05-engel}).

\casehead{Contribution}
The paper introduced a demand system that is flexible, consistent with utility maximization, aggregable
without linear Engel curves, and nearly linear to estimate. Its application established three lasting
practices. The first is to estimate the linear approximation and check it against the exact index. The
second is to test homogeneity equation by equation and symmetry on the system. The third is to check
concavity observation by observation. The application also carries a lesson for industrial organization. A static demand system
estimated on time series can reject the restrictions of theory for reasons that have nothing to do with
preferences, such as dynamics, aggregation or trends. Applications in industrial organization use panels
of markets with market-specific intercepts, partly to absorb such omitted factors. Examples are the
product-by-region effects of \citet{ChaudhuriGoldbergJia2006} and the city effects of \citet{Hausman1997}.

\begin{table}[htbp]
\centering
\caption{Identification and counterfactual map of the British application of Deaton and Muellbauer (1980)}
\label{tab:05-dm-map}
\footnotesize
\renewcommand{\arraystretch}{1.1}
\begin{tabularx}{\textwidth}{@{}>{\raggedright\arraybackslash}p{2.6cm}>{\raggedright\arraybackslash}X@{}}
\toprule
Observed data & Annual British budget shares and price indexes for eight groups of nondurables, with total
  expenditure on the eight groups, 1954--1974. There are 21 observations. \\
\addlinespace
Latent objects & The cost function of a representative household, the aggregation index $\kappa$ and the
  share disturbances. \\
\addlinespace
Estimated parameters & Intercepts $\alpha_i$, expenditure coefficients $\beta_i$ and price coefficients
  $\gamma_{ij}$, with $\alpha_0$ fixed a priori. \\
\addlinespace
Identifying variation & Year-to-year co-movement of budget shares with relative prices and real
  expenditure. Prices and expenditure are treated as exogenous. \\
\addlinespace
Maintained assumptions & A representative household with a constant aggregation index, a static model
  with weakly separable nondurables, and the restricted covariance structure \eqref{eq:05-dm-cov}. \\
\addlinespace
Counterfactual and what it holds fixed & None. The object is a test of homogeneity, symmetry and
  negativity. \\
\addlinespace
Conclusion & Homogeneity is rejected in levels ($LR = 138.7$ with 7 degrees of freedom), and the symmetric
  estimates violate concavity. In first differences homogeneity is not rejected at 5 percent, which points
  to dynamic misspecification. \\
\addlinespace
One limitation & Twenty-one observations support a system with dozens of parameters, so the tests rest on
  asymptotic approximations in a short sample. \\
\addlinespace
Source and our additions & The tables reproduce the article and the authors' book. The exception is Panel
  B of \cref{tab:05-dm-table4}, whose likelihood-ratio statistics and $p$-values are our computation. \\
\addlinespace
Tools in the appendices & \Cref{sec:A-ols,sec:A-mle,sec:B-duality}. \\
\bottomrule
\end{tabularx}
\end{table}

\subsection{Endogenous prices, endogenous expenditure and measured group prices}\label{sec:05-endogeneity}

The British application treats prices and expenditure as exogenous. In the data that industrial
organization uses, brand or product prices in many markets and periods, this assumption usually fails.
This subsection explains why, describes the available instruments, and shows how group prices built from
many stock-keeping units add a measurement problem of their own.

\casehead{Why prices are endogenous}
Firms set prices knowing at least part of the demand shock. A brand-specific advertising success or a
regional epidemic raises $\xi_{it}$ and, through the firm's first-order condition, the price of good $i$.
In the single-product case of \cref{ch:theory}, $p_{it} = mc_{it}\bigl(1 +
\frac{1}{\varepsilon_{ii}}\bigr)^{-1}$, so the price moves with any demand shock that moves marginal cost
or the elasticity. Then $\Cov(\ln p_{it}, \xi_{it}) \neq 0$, and OLS estimates of $\gamma_{ii}$ are biased,
typically toward less elastic demand. Regulated prices, as in Indian pharmaceuticals, weaken this channel
(\cref{sec:05-cgj-identification}).

\casehead{Why group expenditure is endogenous}
At the bottom of a multi-level system the budget variable is group expenditure $X_G$. Within-group share
errors sum to zero, so a pure reallocation inside the group does not change $X_G$. But shocks that raise
demand for the whole group, such as an epidemic that calls for one type of drug, usually raise demand
for particular goods in it as well. The upper-level shock and the lower-level shock
$\xi_{it}$ are then correlated, and so are $\ln\frac{X_{Gt}}{P_{Gt}}$ and $\xi_{it}$.
\citet[p.~219]{Hausman1997} notes that instruments ``might still be required for the segment-expenditure
variable'' even when prices are predetermined. Natural instruments are upper-level expenditure, income
and the price indices of other groups. \citet{ChaudhuriGoldbergJia2006} include total antibiotic
expenditure among their instruments.

\casehead{Instruments for product-space systems}
An unrestricted $J$-good AIDS with homogeneity imposed has $J - 1$ relative prices and real expenditure on
the right-hand side of each share equation. Each equation therefore needs at least $J$ excluded
instruments, and the requirement grows with the number of goods. Three families of instruments are used.
\begin{enumerate}[label=(\alph*)]
  \item \emph{Cost shifters}, such as input prices, are the textbook choice. They are often too few, or
  observed too infrequently, to identify many price coefficients \citep[p.~219]{Hausman1997}.
  \item \emph{Prices of the same product in other markets.} Suppose prices in city $n$ follow $\ln p_{jnt}
  = \rho_j\ln c_{jt} + \mu_{jn} + \zeta_{jnt}$ \citep[eq.~8]{Hausman1997}. Here $c_{jt}$ is a cost
  component common to all cities, $\mu_{jn}$ a city-brand effect, and $\zeta_{jnt}$ a local shock that is
  independent across cities. Prices of brand $j$ in other cities are then correlated with its price in
  city $n$ through $c_{jt}$. They are uncorrelated with the demand error in city $n$. \Cref{sec:05-hlz-identification} develops these instruments and
  \cref{sec:05-bresnahan-critique} presents the critique that national demand shocks invalidate them.
  \citet{Nevo2001} adopts them in characteristic space (\cref{ch:blp}).
  \item \emph{Institutional instruments.} \citet{ChaudhuriGoldbergJia2006} use the number of
  stock-keeping units in each product group, which changes with entry and exit. They also use the prices
  of the five largest units in each group (\cref{sec:05-cgj-identification}).
\end{enumerate}
Characteristic-based instruments, which drive identification in \cref{ch:logit} and \cref{ch:blp}, help
little here. In product space every product has its own intercept and price coefficients, and a product's
characteristics are fixed over time and absorbed by its intercept.

A different route treats prices as predetermined. If retailers fix prices for the week under roughly
constant marginal cost and do not adjust them to that week's demand, prices are uncorrelated with the
week's demand shock. The need for instruments can then be tested with a specification test that compares
the IV and OLS estimates \citep{Hausman1978,Hausman1997}.

\casehead{Estimators}
With instruments $z_t$, the moment conditions are $\E[z_t\otimes\xi_t(\theta)] = 0$, where
$\xi_t(\theta)$ stacks the residuals of the $J - 1$ retained equations of \eqref{eq:05-nl-aids} or
\eqref{eq:05-la-aids}. Equation-by-equation two-stage least squares (2SLS) is consistent but ignores the
cross-equation restrictions. Three-stage least squares, defined in \cref{sec:05-singular}, and GMM impose
symmetry and use the correlation of errors across equations. Two practical points deserve attention. First, consider a two-step estimator that replaces endogenous regressors by their first-stage
fitted values and then runs constrained least squares. It is consistent when the second step is linear in
the replaced variables, since it is then 2SLS. Its second-step standard errors are wrong, however, because
the residual variance is computed with the fitted regressors. Correct 2SLS formulas compute it with the
actual regressors, and the bootstrap is an alternative. \citet[p.~20]{ChaudhuriGoldbergJia2006} use such a two-step procedure. They regress
group prices on the instruments, insert the predicted prices in constrained least squares, and compute
standard errors by the bootstrap (\cref{sec:05-cgj-identification}). Second, when an endogenous variable
enters nonlinearly, plugging in fitted values is
generally inconsistent. The translog index contains products $\ln p_{kt}\ln p_{jt}$. If $\ln p_{kt} =
\widehat{\ln p}_{kt} + v_{kt}$ with first-stage errors $v_{kt}$, then
\begin{equation*}
  \E\bigl[\ln p_{kt}\ln p_{jt}\,\big|\, z_t\bigr] = \widehat{\ln p}_{kt}\,\widehat{\ln p}_{jt}
  + \Cov(v_{kt}, v_{jt}\,|\, z_t),
\end{equation*}
so the plug-in regressor misses a term. The intercepts absorb it only when the conditional covariance is
constant, and even then the intercepts, which also enter the index, are distorted. Nonlinear GMM, which
computes residuals from the actual variables and uses instruments only in the moment conditions, avoids
the problem.

\casehead{Group prices as noisy proxies}
Multi-level systems face a further problem at the bottom level. The price of a product group is an
average of the prices of many stock-keeping units (SKUs), $p_{jt}^A = \sum_{k\in j}\omega_k p_{kt}$ with
revenue weights $\omega_k$. The model is written in terms of an exact group price index $p_{jt}^T$, which
is not observed. Following \citet[pp.~18--19]{ChaudhuriGoldbergJia2006}, let $\varsigma_{it}$ be the
demand shock of the share equation written in the exact indices, and write $\ln p_{jt}^A = \ln p_{jt}^T +
m_{jt}$, where $m_{jt}$ is the proportional approximation error. Let $p_t^T$ and $p_t^A$ collect the exact
and the measured group prices, and let $P(\cdot)$ be the price index of the share equation. Only $P(p_t^A)$
can be computed. Substituting $\ln p_{jt}^T = \ln p_{jt}^A - m_{jt}$ into the share equation, and replacing
$P(p_t^T)$ by $P(p_t^A)$, gives the equation taken to the data,
\begin{equation}\label{eq:05-meas-error}
  w_{it} = \alpha_i + \sum_j\gamma_{ij}\ln p_{jt}^A + \beta_i\ln\frac{X_t}{P(p_t^A)} + \xi_{it},
  \qquad \xi_{it} = \varsigma_{it} - \sum_j\gamma_{ij}m_{jt} + \beta_i\bigl[\ln P(p_t^A) - \ln P(p_t^T)\bigr].
\end{equation}
With $\gamma$ symmetric, for the translog index, to first order in the errors, $\ln P(p_t^A) - \ln P(p_t^T)
\approx \sum_k\bigl(\alpha_k + \sum_j\gamma_{kj}\ln p_{jt}^T\bigr)m_{kt}$. \citet[p.~19]{ChaudhuriGoldbergJia2006} hold the index fixed in
this substitution, so their composite error has only the first two terms. That form is exact only if the
index itself is observed without error, which requires the unobserved $p_{jt}^T$.

The composite error $\xi_{it}$ has two parts, and each is correlated with the observed group prices. The
first, $-\sum_j\gamma_{ij}m_{jt} + \beta_i[\ln P(p_t^A) - \ln P(p_t^T)]$, is correlated with $\ln p_{jt}^A$ by
construction, as in any errors-in-variables problem. The second, the demand shock $\varsigma_{it}$, is
correlated with $\ln p_{jt}^A$ because the weights $\omega_k$ respond to demand.

A one-price version shows how the two channels bias least squares. Suppose $w_t = \alpha + \gamma\ln
p_t^T + \varsigma_t$ and we observe $\ln p_t^A = \ln p_t^T + m_t$, with $m_t$ uncorrelated with $\ln
p_t^T$. The measured price may still be correlated with the demand shock, for example through the
weights. The regression of $w_t$ on $\ln p_t^A$ has error $\varsigma_t - \gamma m_t$. Since
$\Cov(\ln p_t^A, m_t) = \Var(m_t)$ and $\Var(\ln p_t^A) = \Var(\ln p_t^T) + \Var(m_t)$, its probability
limit is
\begin{equation}\label{eq:05-meas-bias}
\begin{split}
  \plim\hat\gamma_{OLS}
  &= \gamma + \frac{\Cov(\ln p_t^A, \varsigma_t - \gamma m_t)}{\Var(\ln p_t^A)} \\
  &= \gamma\,\frac{\Var(\ln p_t^T)}{\Var(\ln p_t^T) + \Var(m_t)}
   + \frac{\Cov(\ln p_t^A, \varsigma_t)}{\Var(\ln p_t^T) + \Var(m_t)} .
\end{split}
\end{equation}
The first term is attenuation toward zero. The second is simultaneity. When positive demand shocks raise
the measured price, it also pushes a negative $\gamma$ toward zero, so both channels make demand look
less elastic. Even if every SKU price were exogenous, OLS on \eqref{eq:05-meas-error} would be
inconsistent, because the measurement error alone biases it. A valid instrument must be correlated with
the observed group price and uncorrelated with both the demand shock and the approximation errors
$m_{jt}$. Chaudhuri, Goldberg and Jia point out that the number of SKUs fails this test if the
approximation error grows with the number of SKUs in a group (their fn.~22).
\Cref{sec:05-cgj-identification} examines their instruments in detail.
